% Options for packages loaded elsewhere
\PassOptionsToPackage{unicode}{hyperref}
\PassOptionsToPackage{hyphens}{url}
\PassOptionsToPackage{dvipsnames,svgnames,x11names}{xcolor}
\documentclass[
  11pt,
]{article}
\usepackage{xcolor}
\usepackage[margin=1in]{geometry}
\usepackage{amsmath,amssymb,amsthm}
\newtheorem{theorem}{Theorem}[section]
\newtheorem{lemma}[theorem]{Lemma}
\newtheorem{corollary}[theorem]{Corollary}
\newtheorem{proposition}[theorem]{Proposition}
\newcommand{\Z}{\mathbb Z}\newcommand{\N}{\mathbb N}
\newcommand{\ETC}{\mathrm{ETC}}\newcommand{\spanw}{\operatorname{span}}
\setcounter{secnumdepth}{3}
\usepackage{iftex}
\ifPDFTeX
  \usepackage[T1]{fontenc}
  \usepackage[utf8]{inputenc}
  \usepackage{textcomp} % provide euro and other symbols
\else % if luatex or xetex
  \usepackage{unicode-math} % this also loads fontspec
  \defaultfontfeatures{Scale=MatchLowercase}
  \defaultfontfeatures[\rmfamily]{Ligatures=TeX,Scale=1}
\fi
\usepackage{lmodern}
\ifPDFTeX\else
  % xetex/luatex font selection
  \setmainfont[]{DejaVu Serif}
  \setmonofont[]{DejaVu Sans Mono}
\fi
% Use upquote if available, for straight quotes in verbatim environments
\IfFileExists{upquote.sty}{\usepackage{upquote}}{}
\IfFileExists{microtype.sty}{% use microtype if available
  \usepackage[]{microtype}
  \UseMicrotypeSet[protrusion]{basicmath} % disable protrusion for tt fonts
}{}
\makeatletter
\@ifundefined{KOMAClassName}{% if non-KOMA class
  \IfFileExists{parskip.sty}{%
    \usepackage{parskip}
  }{% else
    \setlength{\parindent}{0pt}
    \setlength{\parskip}{6pt plus 2pt minus 1pt}}
}{% if KOMA class
  \KOMAoptions{parskip=half}}
\makeatother
\usepackage{longtable,booktabs,array}
\usepackage{calc} % for calculating minipage widths
% Correct order of tables after \paragraph or \subparagraph
\usepackage{etoolbox}
\makeatletter
\patchcmd\longtable{\par}{\if@noskipsec\mbox{}\fi\par}{}{}
\makeatother
% Allow footnotes in longtable head/foot
\IfFileExists{footnotehyper.sty}{\usepackage{footnotehyper}}{\usepackage{footnote}}
\makesavenoteenv{longtable}
\setlength{\emergencystretch}{3em} % prevent overfull lines
\providecommand{\tightlist}{%
  \setlength{\itemsep}{0pt}\setlength{\parskip}{0pt}}
\usepackage{amsmath,amssymb,booktabs}
\setlength{\emergencystretch}{3em}
\usepackage{bookmark}
\IfFileExists{xurl.sty}{\usepackage{xurl}}{} % add URL line breaks if available
\urlstyle{same}
\hypersetup{
  pdftitle={Memory and Spatial Constraints on Eventual Translation in Langton's Ant},
  pdfauthor={Hao Ke},
  colorlinks=true,
  linkcolor={Maroon},
  filecolor={Maroon},
  citecolor={Blue},
  urlcolor={Blue},
  pdfcreator={LaTeX via pandoc}}

\title{Memory and Spatial Constraints on Eventual Translation in Langton's Ant}
\author{Hao Ke}
\date{}

\usepackage{fvextra}
\RecustomVerbatimEnvironment{verbatim}{Verbatim}{breaklines=true,breakanywhere=true,fontsize=\small}
\begin{document}
\maketitle
\begin{abstract}
How much memory and space does a Langton ant need to sustain motion that
never becomes periodically translating? We first give an exact prediction
algorithm. A black read is old at memory length $D$ when its cell was absent
from the preceding $D$ visits. From any complete finite initial state, the next
old-black read is computable, with polynomial bit complexity for fixed $D$.
Finite arithmetic segments retain every skipped write and allow exact continuation.
Geometric bounds give quadratic deadlines for finite event prefixes and total
queries outside any fixed finite region. They imply that a non-eventually-
translating orbit must repeatedly consume black marks that are simultaneously
remote, created late and left unread for a long time. A twenty-cell construction
shows why spatial scale remains necessary in lifetime bounds even when the
intervening event count is fixed. Independently, exact crossing-record counts
and a local visit inequality force whole-prefix oblique width at least
$(1-o(1))\log_2\log T$ by time $T$ on every non-eventually-translating orbit.
Width uses non-axis normals normalized by $\min(|\alpha|,|\beta|)=1$.
Together these results provide effective event queries and necessary temporal
and spatial resources for persistent nontranslation.
\end{abstract}

\section{Introduction}\label{introduction}
A short record of the ant's turns determines the geometry of its recent
visits. If the current cell occurs in that record, its colour follows from
the last visit and the intervening flip. If it does not, a predictor can
provisionally treat it as white. This prediction fails exactly when the ant
reads a black cell absent from its recent visits. Such a cell may be an
unvisited obstacle in the initial colouring or a mark left long ago by the
ant itself.

We study this failure for a memory of $D$ turns. The first question is
algorithmic: given the entire finite colouring now, can one find the next
failure without simulating a potentially distant collision? The second is
physical: after skipping a long predicted flight, can one retain every black
cell needed to continue the calculation? We answer both questions. A finite
suffix predictor gives a repeating candidate motion, exact arithmetic finds
its first disagreement with the colouring, and finite arithmetic segments
record all writes made during the skipped flight.

These algorithms also constrain how old marks can influence later motion.
We prove that a flight ending at a prediction failure lies in a bounded tube
around the segment joining its endpoints. This yields quadratic bounds on
the time to a given number of failures and makes queries outside a fixed
finite region decidable. Applied repeatedly, the bounds imply that an orbit
which never becomes periodically translating must consume black marks that
are simultaneously newly created, arbitrarily old when read, and arbitrarily
far away. Here a mark's lifetime begins with a white-to-black update and ends
at the next read of that cell, which turns it white.

A twenty-cell construction explains the spatial parameter in these bounds. With
only twenty initially black cells and exactly 31 intervening failures for
$D=128$, a family of inputs produces lifetimes of unbounded duration. The
inputs spread farther apart as the duration grows. Thus fixed black mass and
a fixed number of failures do not by themselves control a lifetime.

Crossing records then quantify how much transverse space persistent
nontranslation requires. Equal guarded records force a permanently translating
future. Incoming heights determine the entire crossing history and its retained
inside field. Counting these requests gives a minimum amount of crossing work, while a local corner inequality bounds repeated visits in a narrow
band. Combining the two yields a double-logarithmic lower bound on whole-prefix
oblique width, uniformly over normalized non-axis normals. This argument uses
crossings rather than old-black events, and supplies an independent geometric
constraint on the same class of orbits.

Gajardo's symbolic treatment relates long return paths to infinite memory
effects {[}1{]}. Gajardo, Moreira and Goles establish P-hardness of arbitrary
cell visitation through finite circuit configurations {[}6{]}; their infinite
computation construction uses an infinite background. Our target is the
specific failure of a finite-recency predictor on complete finite input.
For fixed $D$, the event algorithm is polynomial in sparse binary input size.
The periodic read test uses full-drift chronology {[}3{]}, and the finite-region
argument uses the square escape theorem {[}4{]}. Highway multiplicity for
generalized rules is studied in {[}2{]}; we use the classical two-colour rule.

Section~\ref{predicting-the-next-old-black-read} constructs the predictor.
Sections~\ref{sec:endpoint-tube}--\ref{quadratic-event-budgets} prove the
geometric bounds, which then control the cost of exact continuation in
Section~\ref{exact-continuation-with-complete-history}.
Sections~\ref{endpoint-tubes-and-spatially-filtered-feedback}--\ref{sec:birth-feedback}
derive remote feedback. Section~\ref{a-twenty-black-construction-and-the-role-of-distance}
presents the twenty-cell construction and explains the role of spatial scale in the event estimates.
Section~\ref{cross:sec:crossing} derives the independent crossing-information width bound.

\section{Model and results}\label{model-and-results}

A complete state \(X=(B,p,h)\) has finite \(B\subset\mathbb Z^2\). White
turns right, black turns left; flip then move. Cartesian headings N,E,S,W
have vectors $(0,1),(1,0),(0,-1),(-1,0)$ and are numbered $0,1,2,3$.
Write \(\operatorname{pos}(X)=p\), and let \(\operatorname{dir}(h)\) be the unit
vector of heading \(h\). The turn \(\operatorname{turn}(X)\) is
\(\mathrm R\) when \(p\notin B\) and \(\mathrm L\) when \(p\in B\).
The standard highway, abbreviated P104 below, repeats the classical
104-turn word with displacement $(2,2)$ up to quarter-turn rotation and
cyclic phase. Let \(F\) be the
invertible update. A state at time \(t\) precedes update \(t\), which
maps \(F^tX\) to \(F^{t+1}X\). For \(D\ge1\), write \(Q=2^D\) and define the
intrinsic event \[S_D(X)\iff p\in B\quad\text{and}\quad
\operatorname{pos}(F^{-j}X)\ne p\quad(1\le j\le D).\] Its
inverse-history definition applies at the input time as well as later.
It records an old black read. For forward times at least $D$, the definition
uses the actual preceding $D$ visits. At the input time, inverse dynamics
supply the same history convention. An initially black cell can trigger this
event even if the ant has never read it during the forward run.

A \emph{black lifetime} is a pair of update times $b<e$: update $b$ changes a
cell from white to black, and update $e$ is its next read. Its age at consumption
is $e-b$. Initially black inventory has no post-input birth. An orbit is
\emph{eventually translating} if, after some time, its turns repeat with a
heading-restoring period and nonzero displacement; its full black field still
records the accumulated trail. A \emph{complete endpoint} below includes that
entire field, the head position and the heading.

Write $m=|B|$, let $b$ denote the input coordinate bit length when discussing
complexity, and put $R=\max(\{\|q-p\|_\infty:q\in B\}\cup\{0\})$.
The algorithms return either the first requested event with its exact
post-event state, or a certificate that the requested event never occurs.

\textbf{Prediction and continuation.} The next \(S_D\), if present,
occurs before \(Q(R+2Q+10)\), where \(R\) is the initial head-relative
radius. Existence is decidable in polynomial bit time for fixed \(D\) on
sparse explicit input. Finite XOR segments allow the exact full endpoint
to be returned without expanding every skipped write. The first \(n\)
events and their full encoded endpoint are computable in
\(f(D)\operatorname{poly}(m+n+b)\), where \(m\) is the initial mass,
\(b\) the coordinate bit length, and \(n\) is measured numerically.

\textbf{Geometry.} With \(W=\operatorname{diam}_\infty(B\cup\{p\})\),
the time just after event \(n\), whenever it exists, is at most
\begin{equation}\label{eq:event-budget}
T_D(n,W)=Qn(W+2Q+10)+5Q^2n(n-1).
\end{equation} Between a starting head \(p\) and
its next event cell \(e\), the entire finite trajectory stays within
\(3Q+1\) of \([p,e]\). After joining successive event locations the
corresponding tube radius is \(3Q+2\).

\textbf{Remote feedback.} For every fixed finite region \(K\), existence
of an \(S_D\) outside \(K\) is decidable. If a trajectory is not
eventually a nonzero translator, then for every \(D,K,T\) there are
infinitely many black lifetimes born after \(T\), lasting more than
\(D\), and consumed outside \(K\). The properties hold simultaneously
for the same lifetimes.

\textbf{Newly generated feedback.} The first exterior old lifetime born
since the input checkpoint is totally computable, including its exact birth,
consumption and complete post-state. On a non-eventually-translating orbit,
for every sufficiently large time horizon $H$,
Section~\ref{sec:birth-feedback} gives
\(\Omega(\sqrt H/\log H)\) lifetimes with birth after \(H^{1/4}\) up to
integer rounding, radius of order \(\log H\), and age of order
\(\log\log H\).

\textbf{Spatial-scale obstruction.} For every even \(n\ge20\), a
complete 20-black initial state writes a cell at time 60 and next reads
it at \(624n-1516\), with exactly 31 intervening \(S_{128}\) events. Its
first permanent P104 time is \(624n-742\). Thus no lifetime bound
depending only on \(D\), intervening event count and initial mass can
hold uniformly.

\textbf{Crossing information and oblique width.} Independently of prediction
failures, crossing-record counts and a local visit inequality force every
finite-support orbit that is not eventually translating to satisfy
\[
\liminf_{T\to\infty}\frac{w_{\alpha,\beta}(T)}{\log\log T}\ge\frac1{\log2}
\]
for every non-axis normal normalized by $\min(|\alpha|,|\beta|)=1$.
Here $w_{\alpha,\beta}(T)$ is the range of $\alpha x_t+\beta y_t$ over the
whole prefix $0\le t\le T$. Section~\ref{cross:sec:crossing} gives the uniform
finite profile deadline behind this bound.

\section{Predicting the next old-black
read}\label{predicting-the-next-old-black-read}

\subsection{Definitions and an intrinsic finite-local
target}\label{sm-1}

Write \(X\)=(\(B\),\(p\),\(h\)), headings N,E,S,W. Let \(F\) be the literal update. Define

\[
R(X)=\max(\{\|q-p\|_\infty:q\in B\}\cup\{0\}).
\]

For integer $D\ge1$, define the stale-black predicate

\begin{equation}\label{eq:sm1}
S_D(X)\iff p\in B\ \text{and}\ \operatorname{pos}(F^{-k}X)\ne p\quad(1\le k\le D).
\end{equation}

The inverse is literal and unique: set \(q=p-\operatorname{dir}(h)\), flip \(q\), and undo a right
turn when \(q\) is black in the supplied post-state, otherwise undo a left
turn. Only cells at distance at most \(D\) from \(p\) are read during \(D\) inverse
steps. Consequently \eqref{eq:sm1} is a translation-invariant finite local predicate on the present colouring and heading. It is determined by this finite neighbourhood.

Using the intrinsic backward orbit makes \eqref{eq:sm1} defined even at input
time zero. At forward times $t\ge D$, it says exactly: the
current read is black and the same cell has not been read in the
previous \(D\) updates. A very old white reread is not an $S_D$ event. A
never-yet-read initial black cell can be an $S_D$ event too. This
convention is distinguished from an age defined only after time zero.

Let $\tau_D(X)=\min\{t\ge0:S_D(F^tX)\}$, when it exists.
This target comprises infinitely many full states; the result below is
not just specified-full-state reachability.

\subsection{Prediction identity: no event means one unique next
turn}\label{sm-2}

For a \(D\)-letter word \(w=w_0\cdots w_{D-1}\) over
\(\{\mathrm R,\mathrm L\}\), order the letters chronologically from oldest
to newest and execute its geometric turns from a fixed pose. Compare the
final position to the preceding \(D\) pre-update positions, indexed by
\(j=0,\ldots,D-1\). If the last matching position has index \(j\), set
\(f_D(w)=\operatorname{opposite}(w_j)\); if there is no match, set
\(f_D(w)=\mathrm R\). This depends only on \(w\): translation and quarter
turns preserve position equality.

Recover the actual previous \(D\) turns of \(X\) using \(F^{-1}\), and
write \(w(X)\) for the turns at \(F^{-D}X,\ldots,F^{-1}X\), in that order.
Pre-update position \(j\) corresponds to \(F^{-(D-j)}X\), and
\(\operatorname{suffix}_{D-1}(w)=w_1\cdots w_{D-1}\) (the empty word when
\(D=1\)). The following equality is exact for every finite state:

\begin{equation}\label{eq:sm2}
\operatorname{turn}(X)\ne f_D(w(X))\iff S_D(X).
\end{equation}

Moreover the only possible mismatch is predicted \(R\) versus actual \(L\).

Proof: a recent previous read fixes the current colour to its flipped
value. With no recent previous read, $f_D$ predicts \(R\). This is wrong
precisely when the current cell is black. Old white cells and virgin
white cells both give \(R\). The backward suffix is part of the actual
orbit and is recovered uniquely by inverse updates.

While $S_D$ is absent, suffixes evolve autonomously as

\[
w\longmapsto\operatorname{suffix}_{D-1}(w)\,f_D(w).
\]

This is a deterministic map on $Q=2^D$ states. Its generated infinite
turn word has a transient shorter than \(Q\) and period at most \(Q\). This
alone does not prove that the physical ant follows the word forever:
distant black cells may cause a first mismatch. The next section bounds
when any such mismatch can occur.

\subsection{First mismatch against an ultimately periodic motion
plan}\label{sm-3}

Consider a complete finite state \(Y\) of radius \(r\) and a proposed periodic
turn word \(u\) of length \(P\), with matching start/end heading and
displacement $d\ne0$. Let $q_i$, $0\le i<P$, be its relative
pre-update positions in the head-centred frame \(p_Y=0\), so
\(B_Y\) is translated into that frame and \(q_0=0\). Its proposed future reads are

\[
q_i+kd,\qquad k\ge0.
\]

Earlier periodic reads of the same cell correspond exactly to

\begin{equation}\label{eq:sm3}
q_j-q_i=\ell d,\quad \ell\ge0,\quad(\ell>0\ \text{or}\ j<i),\quad k\ge\ell.
\end{equation}

Each such \(\ell\) is at most \(P\): the finite footprint's pairwise
infinity-distance is at most \(P\), and
$\|d\|_\infty\ge1$. Earlier-read
parity can change only at these finitely many thresholds \(\ell\).

The initial colour contribution is $\mathbf1_{B_Y}(q_i+kd)$. If this is
nonzero then

\[
k\|d\|_\infty\le r+\|q_i\|_\infty\le r+P,
\]

so $k\le r+P$. It is a one-epoch contribution, turning on at \(k\) and
off at \(k\)+1. Thus the cell colour produced by literal dynamics at each
phase \(i\), which must match the turn \(u_i\), is

\begin{equation}\label{eq:sm4}
\mathbf1_{B_Y}(q_i+kd)\ \oplus\ \bigl(\#\{\text{eligible pairs in \eqref{eq:sm3}}\}\bmod2\bigr).
\end{equation}

All changes occur by epoch \(r\)+\(P\)+1. Testing every breakpoint, including
the last infinite constant interval, is necessary and sufficient for
infinite execution. If there is a mismatch, its first time is strictly
below

\begin{equation}\label{eq:sm5}
P(r+P+2).
\end{equation}

This statement concerns the first mismatch, not a hypothetical
trajectory after an earlier mismatch. Induction in time ensures that the
proposed path agrees with the real path up to that first point.

If \(u\) does not restore heading, repeat it two or four times to restore
heading. The resulting displacement is zero, since the quarter-rotated
displacement vectors sum to zero. If heading is restored but \(d=0\)
already, use \(u\) itself. In either case let \(v\) be the closed-pose block, of
length at most \(4P\). If two consecutive copies of \(v\) were both physically
valid, all flips would occur twice and the complete state would recur.
Such an orbit would remain in finitely many cells forever, impossible
for the classical ant. Hence a first mismatch occurs before
$2|v|\le8P$ updates.

For clarity, the elementary no-bounded-orbit argument used here is as
follows. If the set I of infinitely read cells were nonempty and finite,
choose a cell maximizing x+y. After the north and east neighbours have
ceased being read, the checkerboard/heading-axis invariant fixes the
only possible entering and leaving directions at that cell. This
requires the same read colour on every sufficiently late visit,
contradicting colour alternation. This proves that the orbit cannot remain
in a bounded region.

\subsection{Uniform waiting theorem}\label{sm-4}

\textbf{Theorem (first-event waiting bound).} For every finite state \(X\) and $D\ge1$, with
$Q=2^D$, put

\begin{equation}\label{eq:sm6}
\mathcal H_D(X)=Q(R(X)+2Q+10).
\end{equation}

Exactly one of the following occurs:

\begin{itemize}
\tightlist
\item
  $\tau_D(X)$ exists and $0\le\tau_D(X)<\mathcal H_D(X)$;
\item
  $S_D$ never occurs at any nonnegative time; the ant has a nonzero
  translating periodic trace with period at most $Q$, beginning at a
  time at most $Q$.
\end{itemize}

Thus observing no $S_D$ event at times 0,\ldots,$\mathcal H_D(X)-1$ proves that
none ever occurs.

Proof: construct the autonomous suffix trajectory. If a physical
mismatch occurs before the first repeated suffix, \eqref{eq:sm2} gives the event
within \(Q\) updates. Otherwise at a repetition time $j\le Q$, the
physical trajectory has followed the plan and a periodic suffix cycle of
length $P\le Q$ begins from its current phase. The current complete
radius is at most \(R\)(\(X\))+\(j\): initial black cells move relative to the ant
by at most \(j\); any newly black cell is also within \(j\) steps of it.

For nonzero translational drift, \eqref{eq:sm5} bounds a first future mismatch by

\[
j+P(R(X)+j+P+2)\le Q(R(X)+2Q+3).
\]

For a closed geometric plan, the other bound is $j+8P\le9Q$. Both
are strictly covered by \eqref{eq:sm6}. If no mismatch exists, the planned
periodic trace is the exact infinite future. A closed or rotating
periodic plan was excluded above, so its drift is nonzero. By \eqref{eq:sm2}, no
mismatch also means no $S_D$ event. This proves both the waiting bound
and the dichotomy.

\subsubsection{What is the corresponding decision
problem?}\label{what-is-the-corresponding-decision-problem}

The family

\begin{equation}\label{eq:sm7}
\mathcal L_D=\{X:\exists t\ge0,\ S_D(F^tX)\}.
\end{equation}

is recursive, uniformly in $D$.

A simple decision algorithm simulates the explicit finite bound \eqref{eq:sm6}.
The implementation uses the stronger symbolic version: only the first
suffix cycle is produced, then all-future ray/flip arithmetic finds the
first event, or verifies a permanent event-free trace. No simulation to
a late arithmetic collision is necessary.

\subsubsection{Bit complexity of the implemented first-event
decision}\label{bit-complexity-of-the-implemented-first-event-decision}

With an explicit sparse list of \(m\) black coordinates of bit length at
most b, the symbolic first-event algorithm uses time $2^{O(D)}$ times
a polynomial in the input length. In particular, for every fixed \(D\),
language \eqref{eq:sm7} is decidable in polynomial time in the sparse
binary-coordinate input size.

Indeed inverse-history recovery uses \(D\) steps; the suffix search uses at
most $Q=2^D$ states; its explicit black field has at most \(m\)+\(Q\)+\(D\) cells.
The candidate word has length at most \(Q\), and the ray verifier checks
$O(Qm+Q^2)$ phase/cell and phase/phase relations. All coordinate,
epoch and offset arithmetic has bit length polynomial in b+\(D\).
Closed-pose control replay costs at most \(8Q\) updates. Ordinary
polynomial-time set/dictionary implementations suffice. The parameter
factor grows exponentially in \(D\).

For fixed \(D\) the upper bound is linear in \(R\). This order cannot in general
be replaced by a radius-independent constant: take a certified
event-free P104 state with \(D\)=128 and insert one remote black cell at a
virgin frontier site \(k\) highway periods ahead. Until its first read the
original periodic trajectory is unchanged. The first event is at
\(104k+O(1)\), while \(R=2k+O(1)\). The phase-ray calculation gives the
constants exactly. The full input coordinates supply the growing scale.

\section{Geometry between prediction failures}\label{sec:endpoint-tube}
A finite event segment starts from a complete state and ends just after its
first $S_D$ update. Its endpoints constrain every intervening read, even
when the segment contains a very long predicted flight. Throughout this
section $Q=2^D$ and $\|\cdot\|$ denotes the maximum norm.

\subsection{Two-endpoint tube for a finite event macro}\label{remote-2}

Start at \(X\) with head \(p\).~Suppose its first $S_D$ event is at offset \(\tau\),
at cell \(e\). The macro includes that read and its actual outgoing move, so
its duration is \(L=\tau+1\). Let \([p,e]\) denote the real line segment
between its endpoints.

\textbf{Theorem (endpoint tube and distance-sensitive time).} Every
head position from time 0 through time \(\tau+1\) is in

\begin{equation}\label{eq:remote3}
[p,e]+[-(3Q+1),3Q+1]^2.
\end{equation}

Every changed cell is in the same tube, and

\begin{equation}\label{eq:remote4}
L\le Q\|e-p\|_\infty+5Q^2.
\end{equation}

The tube does not depend on the radius or cardinality of \(B\). Unchanged
black cells outside it remain physically present; \eqref{eq:remote3} does not
enclose the full input black field. The next-event premise is essential.

\subsubsection{Proof}\label{proof}

Until the first event, the actual turns agree with the autonomous suffix
predictor. Its word has a prefix of length $j\le Q$ and a cycle \(u\) of
length $P\le Q$. If the event occurs in that prefix, its whole path
and actual outgoing move are within \(Q\)+1 of \(p\), sufficient for
\eqref{eq:remote3}--\eqref{eq:remote4}.

Otherwise let a be the head after the j-step prefix, and $v_i$ the
relative pre-update phase positions of \(u\). Then
$\|a-p\|_\infty\le j$ and
$\|v_i\|_\infty\le P$.

First suppose \(u\) restores heading and has nonzero displacement \(d\).~Write
the event as phase i, epoch \(k\), so

\[
e=a+v_i+kd.
\]

When $k\ge1$, any earlier phase position in epoch $\ell\le k$
has

\begin{equation}\label{eq:remote5}
\begin{aligned}a+v_s+\ell d&=(1-\lambda)a+\lambda e+v_s-\lambda v_i,\\ \lambda&=\ell/k\in[0,1].\end{aligned}
\end{equation}

The error has $\ell_\infty$ norm at most \(2P\). Moving a to \(p\) costs at most
\(j\); the prefix is within \(j\) of \(p\) as well. For \(k\)=0 all cycle-prefix
positions are within \(j\)+\(P\) of \(p\) directly. Thus all positions through the
event read lie in $[p,e]+[-3Q,3Q]^2$. The true mismatched
outgoing move costs at most one more.

Also, since \(d\) is a nonzero integer vector,

\[
k\le k\|d\|\le\|e-p\|+j+P.
\]

Hence

\begin{equation}\label{eq:remote6}
\begin{aligned}\tau+1&=j+kP+i+1\\&\le Q\|e-p\|+2Q^2+2Q\\&\le Q\|e-p\|+5Q^2.\end{aligned}
\end{equation}

If the cycle rotates heading, repeat it twice or four times to restore
heading. The resulting geometric block is closed, since its successively
rotated drift vectors sum to zero. If heading is already restored with
\(d=0\), use one copy. The closed block has length at most \(4P\); its footprint
is within \(2P\) of \(a\), because a point on a closed unit-step walk is reached
by either side of the loop. Thus \eqref{eq:remote3} again holds. Two physically
valid copies of a closed block would repeat the complete state (every
cell is toggled an even number of times), impossible by
the no-bounded-orbit argument of Subsection~\ref{sm-3}. Consequently
$\tau+1\le j+8P\le9Q$, which is at most $5Q^2$ because
$Q\ge2$. This proves \eqref{eq:remote4} also in the closed cases. Every
changed cell was a read position, proving the final assertion.
\(\square\)

This argument needs only a finite macro to its first event, not a
presumption that any candidate predictor cycle is physically realizable
forever.

\subsection{The event skeleton carries the large-scale
geometry}\label{remote-3}

List the actual $S_D$ event times $t_1<t_2<\cdots$ and
cells $e_i$. Put

\[
\begin{gathered}H_D=3Q+2,\qquad s_m=t_m+1,\\J_m=\|e_1-p_0\|+\sum_{i=1}^{m-1}\|e_{i+1}-e_i\|.\end{gathered}
\]

\textbf{Corollary (event-chain tube).} Up to the complete post-mth-event state, the
entire head path is contained in the $H_D$-neighbourhood of the polygonal
chain

\begin{equation}\label{eq:remote7}
p_0\longrightarrow e_1\longrightarrow\cdots\longrightarrow e_m.
\end{equation}

and

\begin{equation}\label{eq:remote8}
J_m+1\le s_m\le QJ_m+6Q^2m.
\end{equation}

In particular, all changed cells are in that tube; the complete state
still contains every unaltered initial black cell outside it.

Proof: each later macro starts one unit from $e_i$. Replacing that start
by $e_i$ in \eqref{eq:remote3} adds at most one to the tube radius. Sum \eqref{eq:remote4},
bounding each later source-to-event distance by
$\|e_{i+1}-e_i\|_\infty+1$. The lower bound
follows from unit-speed motion between the actual event cells.
\(\square\)

For every real linear functional a, the maximum of $a\cdot p_s$ over this
prefix lies between $\max(a\cdot p_0,a\cdot e_1,\ldots,a\cdot e_m)$ and that value plus
$H_D\|a\|_1$. Thus, on an infinite event
chain, any direction in which the head is unbounded above is also a
direction in which these event locations are unbounded above. This controls the directional ranges of the event locations.

This is stronger than using the convex hull of the whole black field:
far, untouched black debris does not enlarge the event-to-event tube. The bound follows from the geometry between consecutive events.

\paragraph{A lifetime with no intervening event.}
Let $N_D(b,e)$ count $S_D$ events strictly between the birth and consumption
of a black lifetime. If this count is zero, its age is bounded independently
of the initial diameter. At a
white birth b there is no $S_D$ event. If $e-b>D$ and there is
none in (b,\(e\)), the death is the first $S_D$ from $X_b$ and its
source/event positions are equal. The endpoint-tube time estimate gives
$e-b+1\le5\cdot4^D$. If $e-b\le D$ the weaker bound
$e-b\le5\cdot4^D$ is immediate. Thus

\begin{equation}\label{eq:circuit9}
N_D(b,e)=0\quad\Longrightarrow\quad e-b\le5\cdot4^D.
\end{equation}


\section{Quadratic event budgets}\label{quadratic-event-budgets}

Put $C(X)=\operatorname{conv}(B\cup\{p\})$, and let $X_n$ be the
complete state just after event $n$.

\textbf{Theorem (convex-hull and quadratic event bounds).}
A finite event segment, including every intermediate head position, satisfies
the enclosure
\[C(X_{\rm next})\subset C(X)+[-5Q,5Q]^2.\] Iterating gives \(C(X_n)\subset C(X_0)+[-5Qn,5Qn]^2\) and
the stated \(T_D(n,W)\), together with the radius bound
\(R(X_n)\le W+10Qn\). The
first-event waiting estimate is the predictor bound proved above,
given in \eqref{eq:sm6}.

\subsection{A geometric lemma for a translating motion plan}\label{geo-3}

Start at \(Y=(B,p,h)\). Let a nonempty word \(u\) of length \(P\) restore heading
and have nonzero drift \(d\).~Write its pre-update phase positions as
$q_0,\ldots,q_{P-1}$, with $q_0=p$. The prescribed infinite positions are
$q_i+kd$ for $k\ge0$. These are at first just a motion plan:
executability is not assumed.

Until its first read mismatch with literal LR dynamics, the plan is the
actual orbit, by induction on updates. All earlier plan writes then
legitimately contribute to its current read colour.

\subsubsection{Only finitely many old-field and self-read changes
matter}\label{only-finitely-many-old-field-and-self-read-changes-matter}

A previous phase \(j\) visits the position $q_i+kd$ precisely when

\[
q_j-q_i=sd,\quad s\ge0,\quad(s>0\ \text{or}\ j<i),\quad k\ge s.
\]

For each phase pair this toggles parity at one threshold \(s\). Let $K_{\rm self}$
be the maximum of all these \(s\) and zero. Since the two phase points are
connected by at most \(P\) unit moves within the block,

\begin{equation}\label{eq:geo3}
K_{\rm self}\|d\|_\infty\le P.
\end{equation}

Let $K_B$ be the greatest epoch $k\ge0$ for which some $q_i+kd$
is in the actual finite set \(B\). If there is no such pair, put $K_B$=-1.
This maximum is finite. If $K_B\ge0$, retain one corresponding
black point

\begin{equation}\label{eq:geo4}
z=q_j+K_Bd\in B.
\end{equation}

Put

\begin{equation}\label{eq:geo5}
K_*=\max(0,K_{\rm self},K_B+1).
\end{equation}

At every epoch $k\ge K_*$ the old field contributes zero at every
phase and every self-read threshold has stabilized. The prescribed
read-colour vector is thus constant in \(k\) there. If a mismatch exists at
all, a first mismatch must occur within epochs \(0,\ldots,K_*\), hence at an
offset strictly below \((K_*+1)P\).

This remains true if early old-black cells temporarily mask a
self-inconsistent motion plan. Such masking ends by $K_B$+1 and has not
been silently discarded. The argument works for either direction of
colour mismatch; its later suffix instantiation has the special R-to-L
mismatch property.

\subsubsection{The same cutoff bounds space by a constant halo, not
by travel
time}\label{the-same-cutoff-bounds-space-by-a-constant-halo-not-by-travel-time}

We prove that every prescribed read in epochs \(0,\ldots,K_*\) is in

\begin{equation}\label{eq:geo6}
\operatorname{conv}(B\cup\{p\})+[-3P,3P]^2.
\end{equation}

If \(K_*=K_{\rm self}\), \eqref{eq:geo3} gives
$\|kd\|_\infty\le P$ for all $k\le K_*$,
while $\|q_i-p\|_\infty\le P$. These points
are within \(2P\) of \(p\), already sufficient for \eqref{eq:geo6}. If \(K_*=0\) the statement
is immediate.

Otherwise take \(K_*=K_B+1\) and the old black point \(z\) from \eqref{eq:geo4}. For
$0\le k\le K_B$ and $K_B>0$, put $\lambda=k/K_B$.
Then

\[
q_i+kd=(1-\lambda)p+\lambda z+(q_i-p)-\lambda(q_j-p).
\]

The first two terms lie on the segment \([p,z]\); the remainder has
$\ell_\infty$ norm at most \(2P\). For \(k\)=$K_B$+1, add \(d\) to the \(k\)=$K_B$
expression; $\|d\|_\infty\le P$, giving \(3P\).
When $K_B$=0, the two epochs \(k\)=0,1 are within \(2P\) of \(p\) directly. Thus
\eqref{eq:geo6} holds in every case, including ties between maxima in \eqref{eq:geo5}.

In particular, if the plan has a first physical mismatch, all positions
before and at that read satisfy \eqref{eq:geo6}. Executing the actual mismatched
move instead of the planned one adds at most one unit of halo. The old
field and every new black cell at the resulting complete endpoint also
lie inside \eqref{eq:geo6} enlarged by one, since each new black cell is a
previously updated location.

There is no claim that an arbitrarily long valid event-free translating
tail stays in \eqref{eq:geo6}. The finite mismatch is an essential hypothesis.

\subsection{The one-event convex-hull bound}\label{geo-4}

Run the order-D predictor from \(X\) until an immediate $S_D$ event or a
repeated suffix. Before the repetition $j\le Q$, all written cells
and the head are within \(Q\) of the original head, so

\begin{equation}\label{eq:geo7}
C(X_j)\subset C(X)+[-Q,Q]^2.
\end{equation}

An immediate event and its outgoing move have the claimed bound
trivially. At a repeated suffix the plan cycle has length $P\le Q$.

If the heading is restored and drift is nonzero, apply Subsection
\ref{geo-3} from the full state $X_j$. A finite next event is exactly its
first mismatch. By \eqref{eq:geo6},\eqref{eq:geo7} and the real final update, the entire
macro stays within halo

\begin{equation}\label{eq:geo8}
Q+3P+1\le4Q+1\le5Q.
\end{equation}

If the cycle rotates heading, repeat it at most four times to get a
closed-pose plan; if its displacement is already zero use one copy. This
closed plan has length $L\le4P$. Each point of a closed unit-step
walk is reachable from its origin in either time direction, so its
distance from the origin is at most $L/2\le2P$. If two copies of
that plan were fully executable, they would flip each cell an even total
number of times and restore the pose, repeating the complete state.
Consequently a mismatch occurs before $2L\le8P$. All positions up
to its read are in the same closed footprint, and the true mismatched
move adds one unit. The resulting source-relative halo is at most
$Q+2P+1\le5Q$.

Every unchanged black cell was already in C(\(X\)), and every newly black
cell was on the bounded read footprint. This proves both the path and
complete-support parts of the one-event enclosure.

\subsection{Iteration and coordinate bit lengths}\label{geo-5}

Minkowski sums of $\ell_\infty$ balls add their radii. Applying the one-event
enclosure successively proves the iterated enclosure, and also covers the complete trajectory
between checkpoints. The $\ell_\infty$ diameter of
$C(X_0)+[-5Qn,5Qn]^2$ is $W+10Qn$. Both the current head and each
current black cell lie in that set, proving the stated radius bound.

The nth macro starts after n-1 events. Because the event offset is strictly less than \eqref{eq:sm6}, and the post-event duration is offset+1, its duration is at most

\[
Q(W+10Q(n-1)+2Q+10).
\]

Summing from 1 through \(n\) proves the quadratic time bound. These are integer non-strict
bounds for post-event time; the event itself occurred one update
earlier.

For each fixed $D$, the event deadline is quadratic in the event count.
Its constants retain the initial diameter and depend exponentially on $D$.

\subsubsection{Complete XOR histories still have a growing number of
terms}\label{complete-xor-histories-still-have-a-growing-number-of-terms}

In the symbolic representation of Section~\ref{exact-continuation-with-complete-history}, a skipped periodic phase adds a
finite arithmetic segment of actual read positions, and ordinary updates
add singletons. Starting from an explicit m-black input, every generated
term endpoint up to \(n\) events is an actual read location covered by the iterated enclosure;
original singleton terms also lie in that hull. This is true before
syntactic cancellations, since every term records actual read positions
along the finite trajectory. Each nontrivial segment has
nonzero integer stride, so its length is at most W+10Qn+1 (distance
between its two endpoints is length-1 times a norm at least one). If
initial absolute coordinate bit length is b, endpoints and lengths thus
have bit length

\begin{equation}\label{eq:geo9}
O(b+D+\log(n+1)).
\end{equation}

per coordinate and length. The number of segment terms can grow as
$m+n(9Q+1)$, as proved in the continuation construction below. Temporary inverse-history reconstruction
adds at most \(D\) additional local moves; it changes neither the returned
state nor this asymptotic bit bound.

\section{Exact continuation with complete
history}\label{exact-continuation-with-complete-history}

\subsection{Finite arithmetic segments and symmetric difference}\label{xor-2}

For integer a,\(v\) and $n\ge1$ put

\[
\operatorname{Seg}(a,v,n)=\{a+jv:0\le j<n\}.
\]

If $n>1$ require $v\ne0$. An encoded colouring consists of a
finite list of segments and means their symmetric difference. A
singleton is a segment of length one. Thus an explicit finite black set
embeds exactly as singleton terms. Every encoded state still has finite
support, including after a huge skipped prefix. Overlaps cancel modulo
two; they are never silently interpreted as union.

Membership of a point is the parity of the segments containing it. For a
nontrivial segment, containment is one integer-collinearity test
followed by a finite-index interval test. A literal update and its
inverse therefore remain exactly computable without expanding any
segment. Identical singleton terms may be cancelled syntactically; no
cell is removed because it is presumed irrelevant. The list is not
canonical: distinct expressions may encode the same finite set.

\subsection{Exact first mismatch against a periodic motion
plan}\label{xor-3}

Let a supplied word \(u\) have length \(P\), restore heading, and have nonzero
drift \(d\). Its pre-update phase positions are $q_i$. The proposed future
reads are

\[
q_i+kd,\qquad0\le i<P,\quad k\ge0.
\]

As in the earlier finite-field verifier, an earlier periodic read at
this same position exists for each pair (\(j\),\(r\)) satisfying

\[
q_j-q_i=rd,\quad r\ge0,\quad(r>0\ \text{or}\ j<i),\quad k\ge r.
\]

Each pair toggles read parity at the one threshold \(r\). These are finite
in number.

Now intersect a phase ray with one old finite segment:

\begin{equation}\label{eq:xor1}
q_i+kd=a+jv,\qquad k\ge0,\quad0\le j<n.
\end{equation}

If \(d\) and \(v\) are not parallel, determinants give at most one integer pair
(\(k\),\(j\)), which is tested against the bounds. If they are parallel but
a-$q_i$ is not, there is no intersection. Otherwise reduce a nonzero
coordinate of \(d\) to

\[
\alpha k-\beta j=\gamma,\qquad \alpha>0.
\]

Writing \(g=\gcd(\alpha,|\beta|)\), there is no solution unless \(g\)
divides \(\gamma\). When it does, solutions for \(k\) form one residue class modulo
\(|\beta|/g\). The finite bounds on \(j\) supply an interval
\([\mathrm{lo},\mathrm{hi}]\), intersected with $k\ge0$. Thus \eqref{eq:xor1} is either
empty, a singleton, or a finite arithmetic progression of epochs. The
singleton case also handles length-one terms independently of the unused
stride.

For a phase i, its initial-black contribution is the XOR of membership
in these finitely many epoch progressions. Partition $k\ge0$ at
all progression starts, ends+1, and earlier-read thresholds. Within each
partition interval, the active progressions are fixed; their combined
parity is periodic in \(k\), with period dividing the least common multiple
of their moduli. The earlier- read parity is constant on that interval.
Test a whole residue period, or the shorter finite interval. The final
unbounded interval has no active finite segment and needs one parity
test.

This is an exhaustive finite algorithm, not sampling a few epochs. It
decides whether $u^\omega$ is the exact future; otherwise it finds the
first mismatch. Until the first mismatch the proposed and literal paths
agree by induction, so its position and time are real. The running time
can depend on a large lcm. No polynomial-time claim is made for
arbitrary segment-encoded input strides.

\subsection{Exact algebra of the skipped writes}\label{xor-4}

Suppose the first \(n\) proposed updates have been verified (including \(n\)=0).
For each phase i put

\[
n_i=\max\!\left(0,\left\lfloor\frac{n-1-i}{P}\right\rfloor+1\right).
\]

Its updated cells are exactly $\operatorname{Seg}(q_i,d,n_i)$, omitting phases with
$n_i=0$. Therefore the complete endpoint is

\begin{equation}\label{eq:xor2}
B_n=B_0\mathbin{\triangle}\mathop{\triangle}_{\substack{0\le i<P\\n_i>0}}\operatorname{Seg}(q_i,d,n_i).
\end{equation}

If \(n\)=kP+i, the endpoint position is $q_i+kd$ and its heading is the
recorded heading before phase i. Equation \eqref{eq:xor2} accounts for every
flip, including self-overlap, overlap between phases, and overlap with
every old segment. It adds at most \(P\) terms regardless of \(n\).~The helper
applying \eqref{eq:xor2} does not assert that an arbitrary geometric word is
physically executable; its caller must first establish validity through
the chosen prefix.

At the first mismatch of the recency-D predictor, the planned turn is \(R\)
and the actual turn is \(L\). Toggle that point in \eqref{eq:xor2}, turn left and
move once. This is the exact full input for the next macro, with no
separate history oracle: the needed preceding \(D\) positions are recovered
by the literal inverse.

\subsection{Total single-event macro and finite iteration}\label{xor-5}

For $D\ge1$ put $Q=2^D$. The previous suffix identity says that
no $S_D$ event means one deterministic next turn from the latest D-turn
word. There are \(Q\) words. Execute at most \(Q\) matching turns until an event
or a repeated suffix occurs. In the latter case the prediction has a
period $P\le Q$.

If its geometric period is nontranslating or rotates the heading, two
copies of a closed-pose block of length at most \(4P\) expose a mismatch:
otherwise the complete state would recur and the ant would be
permanently confined. Literal no-bounded-orbit excludes that. Replaying
at most \(8P\) steps handles this case.

For nonzero translation apply Subsection \ref{xor-3}. A valid infinite
word gives a certified permanent event-free tail. Otherwise \eqref{eq:xor2} and
one real black update give the next full post-event state. Every input
is finite and every arithmetic search here is finite, so this single
macro is total without caller resource caps.

If the input expression has \(M\) terms, an event output has at most

\begin{equation}\label{eq:xor3}
M+9Q+1\quad\text{terms}.
\end{equation}

The loose uniform value covers \(Q\) preliminary writes, at most \(8Q\)
closed-plan writes, or at most \(Q\) new periodic segment terms, and the
final event. Inverse history recovery is temporary and does not alter
the returned physical field. Nontrivial new segment strides are period
drifts; both coordinate magnitudes are at most \(Q\). The same statement
applies to a terminating tail's entry state.

After $n$ successive event outputs from an explicit $m$-black input,
the expression has at most $m+n(9Q+1)$ terms. The finite iteration computes
the first $n$ events, or stops earlier with a certified event-free tail.
The returned field includes every initial cell and every subsequent flip.

\subsection{Complexity of a finite event prefix}\label{xor-6}

For an explicit sparse input with \(m\) black cells, coordinate bit length
b, fixed \(D\), and an event count \(n\) measured numerically (or in unary), the
first \(n\) events and their complete compressed endpoint are computable in

\begin{equation}\label{eq:xor4}
f(D)\operatorname{poly}(m+n+b).
\end{equation}

All initial
terms are singletons, and every nontrivial generated stride has
components of magnitude at most \(Q\) by \eqref{eq:xor3}. A parallel intersection
modulus is \(|\beta|/g\le Q\), where \(\beta\) is the chosen
signed coordinate of the segment stride after the same sign normalization
as \(\alpha\) in Section~\ref{xor-3}. Hence all residue periods
divide $\Lambda_Q=\operatorname{lcm}(1,\ldots,Q)$, a function only of \(D\). The numbers of
old terms, phase pairs and interval endpoints are polynomial in \(m\)+\(n\) for
fixed \(Q\). One whole residue enumeration therefore costs a factor at most
$\Lambda_Q$, not a factor equal to the huge skipped physical time.

The convex-hull bound in Section~\ref{quadratic-event-budgets} gives
$O(b+D+\log(n+1))$ bits for every generated endpoint and segment length.
Membership, determinant, gcd and residue computations therefore have
polynomial bit cost. Together with the polynomial number of terms and the
fixed-$D$ residue enumeration, this proves \eqref{eq:xor4}.

The generated-stride premise supplies the fixed residue bound used in
\eqref{eq:xor4}; arbitrary symbolic strides are discussed in
Section~\ref{limitations}.

\section{Spatially filtered events and remote feedback}\label{endpoint-tubes-and-spatially-filtered-feedback}
The tube bounds turn a spatial event restriction into a finite stopping
criterion. We use the following escape theorem from {[}4{]}: for any
axis-aligned square with $\ell$ lattice sites on each side, any colouring
and any starting pose inside it, the first exit occurs within
\[
 E_*(\ell)=\frac{13\cdot3^{\ell-1}-2\ell^2-4\ell-5}{2}
\]
updates. Exterior colours remain unread until the exit. For uniformity with
the supplied event certificates, all bounds below use the valid, looser
choice $E(\ell)=8\ell3^\ell+1\ge E_*(\ell)$, for $\ell\ge1$.

\subsection{A finite-region burst bound, independent of the outside
field}\label{remote-4}

Let \(K\) be a nonempty fixed finite set of lattice cells, \(w\) its $\ell_\infty$
diameter, and choose an axis-aligned square enclosing
$K+[-H_D,H_D]^2$. A sufficient side is

\begin{equation}\label{eq:remote9}
\ell(K,D)=w+2H_D+1,\qquad b(K,D)=E(\ell(K,D)).
\end{equation}

\textbf{Theorem (bounded internal bursts).} A consecutive run in the
event list whose event cells all belong to \(K\) has length at most b(\(K\),\(D\)).
The bound is independent of the rest of the complete black field and of
the starting time.

Proof: start immediately before the first event of that run. The event-chain tube corollary
keeps the whole actual path, through every completed internal event in
the run, inside the square in \eqref{eq:remote9}. Its first exit is no later than
\(E(\ell)\). Thus the elapsed time through any such consecutive completed run
is strictly less than \(E(\ell)\), and the number of distinct event reads is
at most that elapsed time. The slightly weaker integer bound $b(K,D)$ is
sufficient. Exterior colours can be arbitrary because they stay
unchanged before first exit. \(\square\)

The word `consecutive' refers to events, not ordinary updates. The ant
may execute many non-event updates between them. Neither infinite local
oscillation nor arbitrarily much exterior storage invalidates this
uniform bound.

\textbf{Corollary (unbounded event locations).} For each
fixed \(D\), an orbit has infinitely many $S_D$ events if and only if their
positions form an unbounded set. If the events are infinite, infinitely
many lie outside every fixed finite \(K\).

An infinite tail confined to \(K\) would contradict the bounded-burst theorem after its last
external event. A finite event list has a finite position set, giving
the reverse direction. This is not a claim that $e_i$ tends to infinity:
returns to \(K\) can occur infinitely often, interleaved with external
events.

\subsection{A total spatially filtered event
algorithm}\label{remote-5}

\textbf{Theorem (spatially filtered event query).} Uniformly in a complete finite \(X\), integer
$D\ge1$ and explicitly supplied finite \(K\), the following
predicate is decidable:

\begin{equation}\label{eq:remote10}
\exists t\ge0:\quad S_D(F^tX)\ \text{and}\ \operatorname{pos}(F^tX)\notin K.
\end{equation}

The algorithm returns the first such event, when one exists. A negative
answer also certifies eventual nonzero translating periodicity, but need
not certify standard P104.

Algorithm: call the already-total next-event procedure. If it certifies
no future $S_D$, return a negative answer and its periodic certificate.
If its event cell is outside \(K\), return that event and accumulated time.
Otherwise execute the finite prefix through that event with the complete
field retained and repeat.

Every individual query and finite replay terminates. Were the outer loop
infinite, there would be infinitely many $S_D$ events all in \(K\),
contradicting the unbounded-event-location corollary. This proves termination. More
quantitatively, use \(K\) union \{$p_0$\} in \eqref{eq:remote9}: every completed
inside-only prefix has its whole path inside that square, so fewer than
its \(E(\ell)\) literal updates, hence at most \(E(\ell)\) internal events, can
precede the final query. The final external event itself may be
arbitrarily distant in time. For \(K\) empty, the original next-event
procedure suffices directly.

The implementation leaves \(E(\ell)\) symbolic, instead of attempting to
construct an integer such as $3^{6\cdot2^{128}}$. It uses literal replay
of inside-event intervals; no efficient full symbolic continuation is
claimed for huge inside intervals. The symbolic escape bound is used to
prove that the number of iterations is finite.

\subsubsection{An explicit first-external-event
deadline}\label{an-explicit-first-external-event-deadline}

Let $R_0=\max_{q\in B_0}\|q-p_0\|_\infty$
(zero for empty \(B\)), and use $\ell_0=\ell(K\cup\{p_0\},D)$. If an
external event exists, its first time is strictly less than

\begin{equation}\label{eq:remote10a}
(Q+1)E(\ell_0)+Q(R_0+2Q+10).
\end{equation}

Indeed, the head first exits the guard square at some
$\tau\le E(\ell_0)$. Either an external event has already occurred,
or all preceding event cells are in \(K\). In the latter case the next $S_D$
after \(\tau\), if present, cannot be in \(K\): otherwise \eqref{eq:remote7} would enclose
the whole prefix through it inside that square, contradicting the exit
at \(\tau\). Apply the existing first-event waiting bound at $X_\tau$, using
$R(X_\tau)\le R_0+\tau$. If there is no next event, there is no
external event at all. This is also a direct two-stage decision proof,
independent of the loop argument.

For fixed \(D\) and fixed \(K\) in the normalized $p_0$=0 frame, the filtered
predicate is polynomial-time in the sparse binary-coordinate input
encoding, using the existing fixed-D polynomial-time next-event
procedure. There are only constantly many internal events, and the total
literal inside-prefix replay is bounded by $E(\ell_0)$, a constant in this
fixed-parameter assertion.

The spatially filtered query goes beyond the unfiltered first-event test of
Section~\ref{predicting-the-next-old-black-read}: the first many $S_D$
events can occur inside \(K\) and must be ignored. The earlier fixed-n event
theorem alone does not say when that filtering loop stops. The endpoint
tube plus finite-domain escape now proves that it always does.

\subsection{Quantitative exterior feedback on every infinite event
chain}\label{remote-6}

Let $N_D(t)$ count all $S_D$ reads at offsets $0\le s<t$, and let
$M_{D,K}(t)$ count those outside fixed $K$. Put $b_K=b(K,D)$, as defined
in the bounded-burst theorem. Internal events occur in at most
$M_{D,K}(t)+1$ runs, so that theorem gives

\begin{equation}\label{eq:remote11}
\begin{aligned}N_D(t)&\le b_K+(b_K+1)M_{D,K}(t),\\M_{D,K}(t)&\ge\left\lfloor\frac{N_D(t)}{b_K+1}\right\rfloor.\end{aligned}
\end{equation}

For an infinite $S_D$ chain, define the computable integer

\begin{equation}\label{eq:remote12}
n_D(t,W)=\max\{n\ge0:T_D(n,W)\le t\}.
\end{equation}

with the deadline \eqref{eq:event-budget}. It follows that

\begin{equation}\label{eq:remote13}
M_{D,K}(t)\ge\left\lfloor\frac{n_D(t,W)}{b_K+1}\right\rfloor.
\end{equation}

For fixed $X,D,K$, this gives an $\Omega(\sqrt t)$ lower bound on the
number of exterior events along an infinite event chain. The constant
depends on the fixed region through $b_K+1$. For empty $K$, use
$M_{D,K}=N_D$ and $b_K=0$.

\subsubsection{A stronger exterior-only quadratic
deadline}\label{a-stronger-exterior-only-quadratic-deadline}

The factor $b_K+1$ in \eqref{eq:remote13} can be avoided in the leading quadratic
coefficient. Augment the state hull by the observation region:

\[
C_K(X)=\operatorname{conv}(B\cup\{p\}\cup K),\qquad W_K=\operatorname{diam}_\infty C_K(X_0).
\]

A finite macro ending just after the first external $S_D$ event satisfies

\begin{equation}\label{eq:remote13ab}
\begin{aligned}C_K(X_{\rm new})&\subset C_K(X)+[-9Q,9Q]^2,\\\text{duration}&\le2Q\operatorname{diam}_\infty C_K(X)+20Q^2+b_K.\end{aligned}
\end{equation}

To prove \eqref{eq:remote13ab}, if there are internal events, their entire prefix
is in $\operatorname{conv}(\{p\}\cup K)+[-H_D,H_D]^2$ by \eqref{eq:remote7}. The full
field at the last internal post-event state is therefore in
$C_K(X)+[-H_D,H_D]^2$, without dropping the initial field. Apply
the \(5Q\) one-event convex-hull theorem of Section~\ref{quadratic-event-budgets} to the final ordinary
event macro. Since $H_D+5Q=8Q+2\le9Q$, \eqref{eq:remote13ab} follows. With
no internal event, the original \(5Q\) bound is already sufficient.

For the duration calculation, put \(W:=\operatorname{diam}_\infty C_K(X)\).
Then \eqref{eq:remote4} bounds the first internal event duration by
\[QW+5Q^2.\]
From its pre-read state through the last internal post-event state there are
fewer than \(b_K\) further updates, by the burst proof. At that boundary the
complete head-relative radius is at most \(W+2H_D=W+6Q+4\).
The first-event estimate makes the last macro at most \(Q(W+8Q+14)\)
updates. Adding these conservative bounds gives
\[
2QW+13Q^2+14Q+b_K\le2QW+20Q^2+b_K,
\]
since \(Q\ge2\). With no internal event, the original
\(Q(W+2Q+10)\) estimate is smaller. Overcounting the first internal read
only weakens the bound. For empty \(K\), use \(b_K=0\).

After \(n\) external events, \eqref{eq:remote13ab} gives augmented width at most
$W_K+18Qn$. Consequently the post-nth-external-event time, if present,
satisfies

\begin{equation}\label{eq:remote13c}
U_{D,K}(n,W_K)=n(2QW_K+20Q^2+b_K)+18Q^2n(n-1).
\end{equation}

On an infinite chain,

\begin{equation}\label{eq:remote13d}
M_{D,K}(t)\ge\max\{n\ge0:U_{D,K}(n,W_K)\le t\}.
\end{equation}

For every fixed \(K\), the leading asymptotic coefficient is
\(\sqrt{t}/(\sqrt{18}Q)\), independent of \(K\). The spatial expansion
is measured directly in exterior-event count.

\subsection{Simultaneously remote, late-born, arbitrarily long black
lifetimes}\label{remote-7}

An orbit that is not eventually a nonzero translating periodic trace has
infinitely many $S_D$ events for each $D$, by the first-event waiting theorem.

\textbf{Theorem (remote long feedback necessity).} On every such
orbit, for every $D\ge1$, finite \(K\), and time T, there are
infinitely many cells/lifetimes described by times b,\(e\) with

\begin{equation}\label{eq:remote14}
\begin{gathered}T<b<e,\qquad e-b>D,\qquad x\notin K,\\\text{update }b\text{ reads }x\text{ white and makes it black},\\\text{update }e\text{ is its next read, black}.\end{gathered}
\end{equation}

These lifetimes can be chosen at infinitely many distinct, unbounded
cells. This does not assert that the whole event sequence eventually
avoids each cell.

Proof: restart at the complete state at time \(T+1\). The unbounded-event-location corollary supplies
infinitely many later $S_D$ events outside \(K\). At most
$|B_{T+1}|$ can consume black lifetimes already present
at the restart. Every other such read consumes a lifetime born at a real
forward white update $b\ge T+1$, at the same exterior cell. The
intrinsic $S_D$ condition then forces $e-b>D$. There was no
intermediate read because black lifetimes end at their next read.
Removing the finitely many initial-lifetime events from an unbounded
event-position set leaves an unbounded set, so the surviving lifetimes
can be chosen at infinitely many distinct cells. Equations \eqref{eq:remote13}
and \eqref{eq:remote13d} also give count bounds after subtracting the finite
stock. \(\square\)

The same lifetimes satisfy all three requirements: late birth, large age
and remote consumption. The endpoint-tube argument forces this feedback to
recur outside every finite region.
 \section{Exact detection of newly generated remote feedback}
\label{sec:birth-feedback}
Spatial filtering can be refined to distinguish newly generated memory from
black lifetimes already present in the input. This refinement returns the
first qualifying lifetime, its provenance, and a complete state from which
the query can be repeated. It also makes the dependence of remote feedback
on the observation scale quantitative.

\subsection{The first birth-filtered query}
Fix a complete source state \(X=(B,p,h)\), \(D\ge1\), and a finite set
\(K\). A qualifying lifetime consists of consecutive reads of a cell
\(q\notin K\) at times \(b,e\), with
\[
 0\le b<e,\qquad e-b>D,
\]
where update \(b\) turns right and update \(e\) turns left. The question
asks for the least consumption time \(e\), its birth \(b\), the cell
\(q\), and the complete state \(F^{e+1}X\). An initially black coordinate
can qualify after its initial lifetime has been consumed and the site has
been written black again.

Put \(M=|B|\), \(Q=2^D\), and
\(W=\operatorname{diam}_\infty(B\cup\{p\}\cup K)\). To keep one explicit
bound throughout this section, use the conservative square escape estimate
\[
 \ell=\operatorname{diam}_\infty K+2(3Q+2)+1,
 \qquad b_K=8\ell3^\ell+1
\]
for nonempty \(K\), and \(b_K=0\) for empty \(K\). The stronger escape
bound used earlier can decrease this budget. The exterior-event estimate is
\[
 U(n,W)=n(2QW+20Q^2+b_K)+18Q^2n(n-1).
\]

\textbf{Theorem (first newly generated exterior lifetime).}
The birth-filtered query is total on every finite input. It uses at most
\((M+1)(b_K+1)\) ordinary next-\(S_D\) calls. A positive answer satisfies
\(e+1\le U(M+1,W)\). A negative answer certifies eventual nonzero translating
periodicity.

\textit{Proof.}
Process the exact successive \(S_D\) events in chronological order. Ignore
those in \(K\). For an exterior black read, find the last read of the same
cell at a nonnegative time strictly before the event. If one exists, it is
a white read and hence the birth of the consumed lifetime; return it.
If none exists, the event consumes one of the \(M\) source lifetimes.
Each source lifetime can be skipped only once. In particular, consuming a
source lifetime does not permanently exclude its coordinate.

There are at most \(M\) skipped exterior events. Each intervening run of
internal events has length at most \(b_K\), by the fixed-region burst bound.
Either a qualifying exterior event or a permanently event-free suffix is
therefore reached within the stated number of calls. The successful event
is no later than exterior event \(M+1\), which proves the deadline. In the
negative case only finitely many \(S_D\) occur, so the finite-suffix theorem
gives eventual translating periodicity. The ordered search returns the first
qualifying consumption. \(\square\)

The negative certificate concerns translating periodicity; its translating
word is not presumed standard. Optional implementation limits return an
inconclusive status. The unlimited algorithm terminates by the proof above,
without expanding the possibly enormous integer \(3^\ell\).

\subsection{Complete endpoints and exact birth provenance}
Alongside the XOR field, store each actual read sequence as a finite
arithmetic segment with its time progression:
\[
 (a,v,L,t_0,\delta_t,\chi),\qquad
 \text{read }a+jv\text{ at }t_0+j\delta_t,\quad 0\le j<L.
\]
Here \(\chi\in\{\mathrm L,\mathrm R\}\) records the turn at each read in
this phase stream: \(\mathrm L\) reads black and \(\mathrm R\) reads white.
A repeated heading-restoring word contributes one such segment per phase;
literal updates contribute singleton records. These records partition all
executed updates, including those skipped by arithmetic acceleration.
Intersecting a queried cell with every segment and selecting the largest
recorded time below \(e\) recovers its exact last writer. The same records
retain every field toggle, giving the complete reusable post-state.

The arithmetic applies to nonprimitive strides as well. Recall the
phase-ray intersection \(q_i+kd=a+jv\) in \eqref{eq:xor1}, with
\(d\) the period drift, \(q_i\) the phase position, \(k\ge0\), and
\(0\le j<L\). For nontrivial parallel segments write \(d=gu\),
\(v=hu\), \(a-q_i=\eta u\), with \(u\) primitive and \(g>0\).
The intersection equation is \(gk=\eta+hj\); its congruence and finite
index bounds give a finite arithmetic progression of epochs \(k\).
Singleton records use the direct ray-membership test. Nonparallel
segments use an exact determinant solve. Splitting epochs modulo the least
common multiple of progression steps reduces the colour parity to finite
interval sweeps, including the final infinite interval. Thus neither a
remote initial marker nor a remote write is lost during acceleration.

For fixed \(D\) and fixed \(K\) in the translated source-pose frame, the
query has polynomial bit complexity on explicitly listed sparse input.
There are \(O_{D,K}(M+1)\) calls, each adding \(O_D(1)\) segments. Generated
strides and the required residue modulus are bounded in terms of \(D\);
the existing finite-macro bounds control coordinate and time bit lengths.
For arbitrary compressed initial fields, totality still holds, while the
call bound uses actual black cardinality rather than the number of XOR terms.

\subsection{Adaptive checkpoints}
On a trajectory that is never eventually translating periodic, the query
is positive from every actual checkpoint \(s\). One may choose \(D\) and a
finite \(K\) using the known history, then obtain a lifetime born at or after
\(s\), consumed outside \(K\), with age greater than \(D\). Restarting after
its consumption produces temporally nonoverlapping lifetimes. Including
previously selected cells in the next \(K\) makes their sites distinct.
The region is fixed during each call and may change between calls. Each
round has its own complete-source mass and deadline.

\subsection{A simultaneous growing-scale lower bound}
Normalize \(p_0=0\), and put \(M_0=|B_0|\) and
\(W_0=\operatorname{diam}_\infty(B_0\cup\{0\})\). For a large integer
horizon \(H\), define
\[
\begin{gathered}
 L=\lfloor\log_3 H\rfloor\ge64,\qquad
 D_H=\lfloor\log_2 L\rfloor-5,\quad Q_H=2^{D_H},\\
 R_H=\lfloor L/16\rfloor,\quad
 K_H=[-R_H,R_H]^2\cap\mathbb Z^2,\quad
 s_H=\lfloor H^{1/4}\rfloor.
\end{gathered}
\]
Use \(\ell_H=2R_H+6Q_H+5\) and the preceding definition of \(b_K\).
Let \(n_H\) be the largest nonnegative integer satisfying
\[
 U_{D_H,K_H}(n_H,W_0+2R_H)\le H.
\]

\textbf{Theorem (growing-scale feedback).}
On every non-eventually-translating finite-support orbit, the number
\(F(H)\) of black lifetimes satisfying simultaneously
\[
 s_H\le b<e<H,\qquad e-b>D_H,\qquad \|q\|_\infty>R_H
\]
is at least \(\max(0,n_H-M_0-2s_H)\). Moreover,
\[
 \liminf_{H\to\infty}
 \frac{F(H)\lfloor\log_3H\rfloor}{\sqrt H}
 \ge \frac{32}{\sqrt{18}}.
\]

\textit{Proof.}
For each horizon the orbit has infinitely many \(S_{D_H}\), so the uniform
exterior deadline supplies at least \(n_H\) exterior consumptions before
\(H\). At most \(s_H\) precede the cutoff \(s_H\). At that cutoff there
are at most \(M_0+s_H\) black lifetimes, so at most that many later
consumptions can have earlier births. Subtracting these supplies proves
the finite bound. The counted objects are distinct lifetimes, and may occur
at the same coordinate at different times.

For the asymptotic estimate,
\[
 L/64<Q_H\le L/32,\quad
 \ell_H\le5L/16+5,\quad b_{K_H}=O(LH^{5/16}).
\]
The linear coefficient of \(U\) is
\(O_X(L^2+LH^{5/16})\), negligible at its positive quadratic root. Hence
\(n_H=(1+o(1))\sqrt H/(\sqrt{18}Q_H)\). Finally,
\(M_0+2s_H=o(\sqrt H/L)\), giving the displayed lower limit.
\(\square\)

Thus one uniform estimate forces late births, a growing excluded radius, and
a growing memory age on the same collection of lifetimes. The conclusion is
a necessary condition on non-eventually-translating motion. It supplies
\(\Omega_X(\sqrt H/\log H)\) lifetimes, while radius, age and birth cutoff
scale respectively as \(\log H\), \(\log\log H\) and \(H^{1/4}\).

The dedicated controls compare complete endpoints and exact birth times
against literal dynamics on 120 inputs, include source-coordinate rebirth,
and exercise arithmetic at epoch \(10^{60}\). On the blank orbit, repeated
queries with \(D=128\), \(K=\varnothing\) return 31 positive rounds before
a periodic negative answer at checkpoint 9990. This illustrates how exact
renewal detection fits an orbit that eventually reaches the standard highway.


\section{A twenty-black construction and the role of
distance}\label{a-twenty-black-construction-and-the-role-of-distance}

The following construction realizes the spatial-scale obstruction with a
fixed number of input cells and a fixed number of intervening events.
Its lifetime has birth \(b=60\), consumption \(e=624n-1516\), and age
\(e-b=624n-1576\). The exact first permanent standard time is
\(e+774\). All twenty cells are specified at input time. We give the
full affine construction and the guard argument for every even
\(n\ge20\). The outgoing port from the four-cell theorem {[}5{]} supplies
the five building blocks used to control this lifetime and its event count.

\subsection{Exact seed and the inherited reference}\label{circuit-2}

The literal rule is white/right, black/left, toggle then move, on
$\mathbb Z^2$. A state at time t precedes update t. Headings are
$\mathrm N,\mathrm E,\mathrm S,\mathrm W=0,1,2,3$. All unspecified input cells are white.

Use the already certified outgoing four-cell port

\[
G=\{(0,-1),(1,1),(1,3),(2,2)\}.
\]

at the origin facing \(\mathrm E\). Let \(\rho(x,y)=(y,-x)\) be clockwise
quarter rotation. At a pose
\((p,h)\), write $G[p,h]=p+\rho^{h-1}G$. For even $n\ge20$ the
complete input black set is

\[
B_n=\bigcup_{i=0}^4G[p_i,h_i].
\]

with the following five poses:

\begin{longtable}[]{@{}lll@{}}
\toprule\noalign{}
$i$ & $p_i$ & $h_i$ \\
\midrule\noalign{}
\endhead
\bottomrule\noalign{}
\endlastfoot
0 & $(0,0)$ & $\mathrm E$ \\
1 & $(n-1,n)$ & $\mathrm N$ \\
2 & $(-n,3n-4)$ & $\mathrm W$ \\
3 & $(-5n+5,-n)$ & $\mathrm S$ \\
4 & $(-3n,-3n+8)$ & $\mathrm E$ \\
\end{longtable}

Start the one physical ant at \((0,0),\mathrm E\). Subsection \ref{circuit-3} proves
that these five four-cell sets are pairwise disjoint, so the input has
exactly twenty black sites.

\subsubsection{Reference data, checked rather than inferred from a long
run}\label{reference-data-checked-rather-than-inferred-from-a-long-run}

Define \(A\) as the first 182 turns of the literal orbit from
\((G,(0,0),\mathrm E)\), and \(U\) as its next 104 turns.
The prefix \(A\) reaches \((5,1),\mathrm E\) with 36 black sites;
Theorem~2 and Section~3 of {[}5{]} certify that the subsequent word is
\(U^\omega\), with drift \((2,2)\). Define \(q_i\) and \(\theta_i\) as the
absolute reference pose at time \(182+i\). The finite phase data used below
are

\[
\begin{gathered}q_{14}=(8,2),\quad \theta_{14}=\mathrm E,\quad U_{14}=\mathrm R;\\q_{61}=(7,8),\quad \theta_{61}=\mathrm N;\\q_{71}=(6,9),\quad \theta_{71}=\mathrm N.\end{gathered}
\]

The reference reads \((0,2)\) white at update 60, does not read it again in
the 182-step prelude, and never reads it on the tail (all tail
\(x\)-coordinates are at least four). It is therefore a newly born,
permanent black lifetime on the unperturbed reference. The constructed
circuit will read it again.

The four initial \(G\) sites are first read at offsets 1,22,32,58. The only
later $S_{128}$ reads of this reference are at offsets 157,280,281, with
respective previous reads 9,144,145. The no-later-events claim is
checked for all future: for each \(\mathrm L\) phase of \(U\), there is a same-cell
preceding phase at distance at most 128 already available by tail
epoch one. The finite prelude/first epoch is checked literally. The
stationary maximum black reread gap on \(U^\omega\) is 116.

The reference's complete 36-cell state passes both the all-epoch
breakpoint parity verifier and the independent
stable-class/clean-seed-and-ray verifier. The latter criterion uses
Jillhewar's supplied-word realization criterion {[}3{]}.
The four-cell port is the construction of {[}5{]}.

\subsection{Universal autonomous execution of the five
stages}\label{circuit-3}

Put \(n=2k\), $k\ge10$. The five claimed stage words, in their
respective rotated poses, are

\begin{equation}\label{eq:circuit2}
\begin{gathered}AU^{k-4}U[0:61],\\AU^{2k-5}U[0:71],\\AU^{4k-6}U[0:61],\\AU^{2k-7}U[0:71],\\AU^{3k-4}U[0:14].\end{gathered}
\end{equation}

The endpoint identities from the phase data give the start times

\begin{equation}\label{eq:circuit3}
\begin{gathered}t_0=0,\quad t_1=52n-173,\quad t_2=156n-440,\\t_3=364n-821,\quad t_4=468n-1296,\quad e=624n-1516.\end{gathered}
\end{equation}

Stages 0 through 3 end at the next table pose. Stage 4 ends at \((0,2),\mathrm E\),
before its next update. This final update is not included in \eqref{eq:circuit2}.

Endpoint algebra alone would not prove these words executable: previous
wakes and future initial ports could cause an earlier divergence. The
next certificate is the missing physical guard.

\subsubsection{Pairwise separation of all stage read
sets}\label{pairwise-separation-of-all-stage-read-sets}

Let $V_i(k)$ be the full set of pre-update read positions in stage i of
\eqref{eq:circuit2}, not its black set and not just its net flips. The following
ten strict separations cover every pair. A row states $\max\ell(V_A)<\beta\le\min\ell(V_B)$.

\begin{longtable}[]{@{}llll@{}}
\toprule\noalign{}
A & \(B\) & $\ell(x,y)$ & $\beta(k)$ \\
\midrule\noalign{}
\endhead
\bottomrule\noalign{}
\endlastfoot
0 & 1 & $y$ & $2k$ \\
2 & 0 & $x$ & $0$ \\
3 & 0 & $y$ & $-2$ \\
4 & 0 & $x$ & $0$ \\
2 & 1 & $x$ & $-2k+1$ \\
3 & 1 & $y$ & $2k$ \\
4 & 1 & $y$ & $2k$ \\
3 & 2 & $y$ & $-2k+1$ \\
2 & 4 & $x-y$ & $-11$ \\
3 & 4 & $x$ & $-6k$ \\
\end{longtable}

Here is the all-parameter checking argument. Each read set is exactly a
union of 182 affine points, 104 affine arithmetic segments, and the
stated 14/61/71 partial-epoch points. For stage \(i\), put
\(r_i\equiv h_i-1\pmod4\) and \(d_i=\rho^{r_i}(2,2)\). Here \(h_i\) is the
heading of the stage pose in the five-port table. Define the full-period
exponents in \eqref{eq:circuit2} by
\[
(R_0(k),\ldots,R_4(k))=(k-4,2k-5,4k-6,2k-7,3k-4).
\]
A repeated-phase segment has the form

\[
p_i(k)+\rho^{r_i}q_j+ad_i,\qquad0\le a<R_i(k),
\]

Its extrema under \(\ell\) are at \(a=0\) or \(a=R_i(k)-1\), according to the sign of
$\ell(d_i)$. They are affine integer functions of \(k\). Partial-epoch positions
use \(a=R_i(k)\). Thus each table row reduces to finitely many inequalities
$\alpha k+\gamma\ge0$. Such an inequality holds for all integers
$k\ge10$ iff $\alpha\ge0$ and
$10\alpha+\gamma\ge0$. These two conditions are checked exactly,
for every point and segment; no sampling or floating-point optimization
is involved.

The certificate mapped in Subsection~\ref{finite-premises} contains 6,832 such inequalities and has
minimum integer slack zero. Every $R_i(k)$ is also checked to be positive
throughout the domain. In particular the table establishes $V_i(k)\cap V_j(k)=\varnothing$. Every initial port cell belongs to the
corresponding reference's first 59 reads, hence to its $V_i$. The twenty
input cells are consequently distinct, and a port belonging to another
stage is outside the whole current read set.

\subsubsection{Why the disjoint read sets
suffice}\label{why-the-disjoint-read-sets-suffice}

At the start of stage i, previous stages have changed only cells in
their $V_j$. Every other initial port is also outside $V_i$. On $V_i$ the
complete true field therefore agrees with a fresh rotated \(G\), with white
everywhere else on that read set. First-difference induction makes the
entire stage word in \eqref{eq:circuit2} literal. Induction over the five stages
proves \eqref{eq:circuit2}--\eqref{eq:circuit3} for the one complete input $B_n$. All
previous black residue is still present outside the current stage reads.

Formula $B_n$ places all gates at time zero. The read-set separation
therefore proves that each distant gate leaves all earlier stages unchanged.

\subsubsection{The actual long-memory return and the count
31}\label{the-actual-long-memory-return-and-the-count-31}

The target \(x=(0,2)\) is not initially black. Stage 0 writes it black at
\(b=60\). It is not read again in that stage, and the read-set separations
exclude it from stages 1--4 up to time \(e\). At time \(e\) the ant reaches
\(x,\mathrm E\) and reads it black; its next turn is \(\mathrm L\), while
the fresh stage-4 reference would require \(\mathrm R\). This is an exact first returned lifetime, not
a geometric near-return.

In \((60,e)\), stage 0 contributes its three offsets 157,280,281. Each of
stages 1--4 contributes all seven reference offsets
1,22,32,58,157,280,281. No other stage has read any of its cells. Its
four initial black first reads occur at absolute times greater than 128,
and so are genuine intrinsic $S_{128}$ events. All other gaps are those of
the reference. Therefore the count is exactly

\begin{equation}\label{eq:circuit4}
3+4\cdot7=31.
\end{equation}

The consumption at time \(e\) itself is not included. This convention
matters.

\subsection{Full-history P104 recovery
certificate}\label{circuit-4}

For the future after the return, retain the complete endpoint $B_e(n)$. It has all five old wakes. We do not
replace the physical input by a cleaned configuration.

Use the fixed unbounded region

\begin{equation}\label{eq:circuit5}
\Omega=\{(x,y):x\ge-16,\ -14\le x+y\le12\}.
\end{equation}

We prove that the restriction of $B_e(n)$ to \(\Omega\) is the same 90-cell
set \(C\) for every $k\ge10$. This is an equality of restrictions
for every parameter in this constructed family.

\subsubsection{All-parameter equality of the protected
field}\label{all-parameter-equality-of-the-protected-field}

The complete endpoint is the XOR of all twenty initial sites and every
actual read position in \eqref{eq:circuit2}, counted with its visit parity. The
stage-1 read set has $x+y>12$. Stage 3 has $x+y<-14$.
Stage 2 has $x-y<-44$, whereas \(\Omega\) entails
$x-y=2x-(x+y)\ge-44$. Those three stages and their initial ports
miss \(\Omega\).

Stage 0 has a fixed prefix and NE period streams. Every epoch
$j\ge6$ has $x+y>12$; all $k\ge10$ include the
first six epochs. Its terminal partial epoch is outside \(\Omega\). Thus its
restriction is constant.

For stage 4, let \(r=3k-4\) and change epoch variable from \(j\) to \(v=j-r\). Every
tail point becomes

\begin{equation}\label{eq:circuit6}
p_4(k)+q_i+j(2,2)=q_i+(-8,0)+v(2,2).
\end{equation}

independent of \(k\). Complete epochs have $-r\le v\le-1$; the
partial epoch has \(v=0\) and $i<14$. At \(k=10\), \(r=26\). Every additional
epoch $v\le-27$ lies below \(x+y=-14\). The stage-4 initial port and
182-step prefix also lie below that line for all $k\ge10$.

Consequently \(C\) is the following finite XOR, restricted to \(\Omega\):

\begin{equation}\label{eq:circuit7}
\begin{gathered}G;\quad\text{all 182 stage-zero prefix reads};\\q_i+j(2,2),\quad0\le i<104,\ 0\le j<6;\\q_i+(-8,0)+v(2,2),\quad0\le i<104,\ -26\le v<0;\\q_i+(-8,0),\quad0\le i<14.\end{gathered}
\end{equation}

The checker identified in Subsection~\ref{finite-premises} evaluates
the finite XOR expression \eqref{eq:circuit7} directly. It obtains 90 cells. It additionally verifies all 1,247 affine
and 208 constant inequalities used above. The genuine full endpoints for
eight parameters are checked against an independent literal replay as a
control, not as the universal argument.

\subsubsection{Protected scattering and exact first standard
time}\label{protected-scattering-and-exact-first-standard-time}

From the reference state \((C,(0,2),\mathrm E)\), execute 774 literal updates. All
positions remain in \(\Omega\). The complete reference endpoint has 162 black
cells and pose \((5,-3),\mathrm E\). Its next 104-word is the cyclic phase-103 shift
of \(U\), with drift \((2,-2)\). Both independent infinite checks
certify its entire future. Every tail phase lies in \(\Omega\), and adding
\((2,-2)\) preserves \(x+y\) and increases \(x\). Thus the whole infinite reference
future stays in \(\Omega\).

The true state at time \(e\) agrees with this reference on all of \(\Omega\). By
first-difference coupling its whole future read trace is the same, for
every \(n\), irrespective of the physically retained exterior wakes. This
proves standardness from \(e+774\).

The reference turn at offset 773 is \(\mathrm R\); the preceding letter in the
backwards 104-periodic extension of the certified tail is \(\mathrm L\). Therefore a
standard future cannot start at or before \(e+773\): any such 104-periodic
word would agree with the known tail over a whole period and hence at
the preceding offset as well. The mismatch contradicts that. Together
with the positive certificate, this proves the exact first time
claimed above, including the exclusion of starts before \(e\). It does not rely
on a finite observed repetition as an infinite certificate.

Every member of the constructed family therefore reaches the standard
highway at the exact time \(624n-742\).

\subsection{Why a lifetime bound needs spatial information}\label{circuit-5}

Let $N_D(b,e)$ count $S_D$ reads strictly between a black lifetime's birth
and its consumption. On all complete finite-support classical inputs
there is no finite-valued function \(F(D,m,M)\), computable or otherwise,
satisfying

\begin{equation}\label{eq:circuit8}
e-b\le F(D,N_D(b,e),|B_0|).
\end{equation}

for every post-input-created consumed black lifetime. In this family the
three arguments are fixed at \((128,31,20)\), while \(e-b=624n-1576\) is
unbounded. The same examples refute a bound in terms of an upper budget
\(m=31\) and \(M=20\). Their excursion radius is unbounded as well.

The no-intervening-event bound \eqref{eq:circuit9} shows that zero events do give a
diameter-independent lifetime bound. The twenty-cell family shows why this
cannot be extended to an arbitrary fixed positive event count.

The quadratic event deadline retains the input diameter and therefore
allows these delays: in this family $|B_0|=20$ while
$\operatorname{diam}(B_0)=\Theta(n)$. The construction varies the input with
$n$ and every member eventually reaches the standard highway.

\subsubsection{A corresponding count of ten-turn violations}\label{the-same-obstruction-for-the-concurrent-j10-event}

Let $\mathcal L_{10}$ contain the cyclic ten-turn factors of the
standard word $U$, and let $J_{10}$ test whether the true ten-turn window
ending now is outside this set. There are
84 allowed factors. On every $X_n$ the number of \(J_{10}\) events strictly in
\((60,e)\) is exactly 264.

To check this uniformly, replace each $U^r$ in \eqref{eq:circuit2} by one \(U\). All
\(r\) are positive. Every removed complete standard cycle contains only
allowed factors; both length-nine boundary contexts are unchanged and
one retained cycle is longer than their combined lengths. Therefore no
forbidden-factor occurrence is lost or created at a splice. Direct
counting on the fixed compressed word gives 264. The same calculation
through the protected recovery and nine following standard turns gives
532 J10 events from absolute update 9 onwards, with the last at \(e+782\).
Later factors are standard. Eight uncompressed physical runs corroborate
the interior count.

Thus initial black mass and the number of intervening ten-turn
violations also leave these lifetime durations unbounded.

\subsection{Exact finite premises and immutable supplement}\label{finite-premises}

The finite premises in this section are regenerated by the twenty-cell
checker \nolinkurl{memory_feedback/verification/scalable_feedback_circuit/check.py}.
Its source, helper \nolinkurl{circuit.py}, and retained \nolinkurl{report.json}
are fixed in the scientific supplement at commit
\href{https://github.com/kehao95/langtons-ant-core/tree/e93252230ea9109908c23a009824d72d30fb91e2/memory_feedback/verification/scalable_feedback_circuit}{\nolinkurl{e93252230ea9109908c23a009824d72d30fb91e2}}.
Paths below are relative to that directory. The checker uses Python 3.10+
with assertions enabled and only the standard library. Copy that directory to a disposable location
and run \nolinkurl{python3 check.py}; compare the resulting JSON with the
immutable retained report. The dependency map is as follows.

\begin{enumerate}
\item \nolinkurl{circuit.reference()} defines \(A,U,q_i,\theta_i\) by literal
replay from \(G\). \nolinkurl{check.reference_controls()} checks the
36-cell endpoint, offsets \(14,61,71\), birth at update 60, seven event
offsets, and the all-future predecessor bound with maximum stationary
gap 116. Its two infinite-tail checks appear in the report under
\nolinkurl{reference.period_parity} and \nolinkurl{reference.ordered_classes};
the gap and offsets are recorded in the same \nolinkurl{reference} object.
The previous reads \(9,144,145\) are obtained by replaying the first
390 updates, as in that function's literal control.
\item \nolinkurl{check.affine_endpoint_controls()} checks the 20 affine
pose/time equalities in \eqref{eq:circuit3}, recorded under
\nolinkurl{endpoints}. \nolinkurl{circuit.separation_certificate()} checks
all ten rows of the separation table, all 6,832 affine inequalities,
their minimum slack zero and positivity of every repeat count; its
output is \nolinkurl{separation}. These are universal affine checks on
\(k\ge10\), not a sampled-parameter inference.
\item \nolinkurl{circuit.protected_core()} evaluates exactly the XOR in
\eqref{eq:circuit7}. \nolinkurl{circuit.protected_region_certificate()}
checks the 1,247 affine and 208 constant inequalities that justify this
restriction for all \(k\ge10\). The \nolinkurl{protected_region} object
records these counts and the 90-cell result; \nolinkurl{recovery.core_sha256}
identifies the sorted coordinate list evaluated by the checker.
\item \nolinkurl{check.recovery_controls()} checks all 774 protected
updates, the 162-cell endpoint \(((5,-3),\mathrm E)\), the phase-103
shift, drift \((2,-2)\), containment of every prefix/tail phase in
\(\Omega\), and the incompatible turn at offset 773. The
\nolinkurl{recovery} object records these results. Its all-future checks
are \nolinkurl{circuit.verify_period()} (parity breakpoints) and
\nolinkurl{check.ordered_class_check()} (full-drift chronology and
clean-seed/ray equality). Region invariance then protects the whole
tail, as proved above.
\item \nolinkurl{check.finite_family_controls()} independently replays
the eight listed parameters, checks full endpoint fields and the exact
31-event count, and records them under \nolinkurl{finite_family}.
The universal count follows from the reference offsets and stage
separation, rather than from those eight controls.
\nolinkurl{check.factor_controls()} records the exact ten-factor counts
\(84,264,532\) and final offset 782 under \nolinkurl{factor_counts}.
\end{enumerate}

The report also records SHA256 digests of both verifier sources. The
trusted finite basis is this explicit integer replay and the two
implemented infinite-tail criteria, together with the analytic coverage
and coupling arguments above; it has not been formalized in a proof
assistant.

\section{Crossing information and spatial growth}\label{cross:sec:crossing}
Remote feedback is a temporal requirement on persistent nontranslation. Crossing
records give an additional spatial requirement. They describe the information
carried by the head through a cut, a method introduced for Turing machines by
Hennie [7]. The following argument proves the needed half-plane coupling directly.
Its local incoming-axis invariant is standard ant geometry [2].
The log-log work-space threshold reviewed in [8] concerns input length; here
the independent variable is the elapsed time of one complete ant orbit.

Write $\ETC$ for eventual translation of position, heading and turn with a
nonzero drift and a fixed period. The complete black field may retain a growing
wake. This is the same translating alternative used in the feedback theorems.

For real $\alpha,\beta\ne0$, normalize
$\min(|\alpha|,|\beta|)=1$ and define
\[
 w_{\alpha,\beta}(T)=
 \max_{0\le t\le T}(\alpha p_t.x+\beta p_t.y)
 -\min_{0\le t\le T}(\alpha p_t.x+\beta p_t.y).
\]
Here $p_t.x$ and $p_t.y$ denote the horizontal and vertical coordinates of $p_t$.
This is normal-coordinate width, not Euclidean width when the normal has a
large norm. Both components must be nonzero.

\begin{theorem}[Width gap]\label{cross:thm:gap}
Every classical finite-support orbit that is not $\ETC$ satisfies, for every
such normal,
\[
 \liminf_{T\to\infty}\frac{w_{\alpha,\beta}(T)}{\log\log T}
 \ge\frac1{\log2}.
\]
The finite bound from which this follows is uniform in the normalized normal.
Consequently the conclusion also holds for time-dependent normal choices,
and for the infimum of these widths over all admissible normals at each time.
\end{theorem}
\subsection{A finite certificate with its own future protection}
For an integer $a$ write $R_a=\{(x,y):x>a\}$. Suppose two times $s<t$ and cuts
$a<b$ satisfy
\begin{align}
 p_t&=p_s+d,\quad h_t=h_s,\quad d=(b-a,\delta),\label{cross:eq:poses}\\
 (B_s\cap R_a)+d&=B_t\cap R_b,\label{cross:eq:field}\\
 p_u&\in R_a\quad(s\le u\le t).\label{cross:eq:guard}
\end{align}
The comparison in \eqref{cross:eq:field} is the complete inside field.

\begin{lemma}[Nested-domain pumping]\label{cross:lem:nested}
Conditions \eqref{cross:eq:poses}--\eqref{cross:eq:guard} imply that the trace from $s$ is
permanently translating with period $t-s$ and drift $d$.
\end{lemma}
\begin{proof}
Until its first exit, the trajectory in a half-plane depends only on its inside
field and starting pose. Thus the two inside states have the same first-exit
duration $\tau\in\N\cup\{\infty\}$. Put $P=t-s$. If $\tau\le P$, the finite
guard fails. If $P<\tau<\infty$, the orbit must exit the smaller $R_b$ between
$t$ and its exit from $R_a$ at $s+\tau$. But the second inside state predicts
its first exit at $t+\tau>s+\tau$. Therefore $\tau=\infty$. Determinism now
identifies the entire two futures up to translation. Since $b>a$, the drift
is nonzero. No exterior cell has been removed from the physical state.
\end{proof}

Set
\[
 a_0=\max\bigl(\{p_0.x\}\cup\{q.x:q\in B_0\}\bigr).
\]
For a finite actual prefix through $T$, the crossing record of the cut
$x=c+\tfrac12$ lists every crossing in chronological order, with its direction
and height. Normalize heights by subtracting their minimum. If $c\ge a_0$ and
$p_T.x>c$, the directions are necessarily $E,W,\ldots,E$.

\begin{lemma}[Complete crossing records]\label{cross:lem:records}
Two equal nonempty normalized records at cuts $a_0\le a<b<p_T.x$ produce a
certificate satisfying Lemma~\ref{cross:lem:nested} at times no later than $T$.
\end{lemma}
\begin{proof}
Both right half-planes are initially white. Their first complete right-side
excursions have translated entry poses and therefore translated read/write
histories until the first exit. Between excursions, left-side updates do not
change the right-side field. Induction over the equal records gives translated
complete inside fields at their last eastward entries, at times $s$ and $t$.
The translation in height is the difference of record minima. The last entry
at $b$ must follow the last entry at $a$, since otherwise ending beyond $b$
would require a later $b$-entry. There is no exit from $R_a$ between $s$ and
$t$, because $s$ is its last entry before $T$. All three conditions follow.
\end{proof}
The records are finite histories, not a query to the infinite future. A far
initial black cell to the right invalidates the hypothesis even when it has
not yet altered the observed finite path.

\subsection{Exact record capacity and crossing work}
Number headings $N,E,S,W$ by $0,1,2,3$. The value $h-x-y\pmod2$ is invariant.
An east crossing of cut $c+\tfrac12$ arrives in column $c+1$; a west crossing
arrives in column $c$. Both arrivals have horizontal heading, so their height
parities differ. Consecutive heights in a crossing record therefore alternate
parity, also after normalization.

For $H\ge0$ put
\[
 E=\lfloor H/2\rfloor+1,\qquad O=\lceil H/2\rceil,
 \qquad U=EO,\qquad V=(E-1)O.
\]
\begin{proposition}[Parity capacity]\label{cross:prop:count}
The number of normalized height words of length $2j+1$, with span at most $H$
and alternating parity, is
\[
 c_H(j)=(H+1)U^j-HV^j,\qquad j\ge0,
\]
where $0^0=1$.
\end{proposition}
\begin{proof}
Over $\{0,\ldots,H\}$ the count before requiring minimum zero is
$E^{j+1}O^j+O^{j+1}E^j=(H+1)U^j$. Words omitting zero use $E-1$ even symbols
and $O$ odd symbols, giving $HV^j$. Subtract them.
\end{proof}
Not every word counted by this proposition is necessarily physically
realizable. It is an exact count of a necessary combinatorial class.

\paragraph{Incoming requests reconstruct the record.}
For an eligible cut the genuine record has the form
\[
 E(y_0),W(z_0),E(y_1),W(z_1),\ldots,E(y_j).
\]
The exit heights $z_i$ are determined responses to the incoming heights $y_i$.
This causal constraint gives a smaller capacity than the parity count.

\begin{lemma}[Request injectivity]\label{cross:lem:requests}
The incoming heights $(y_0,\ldots,y_j)$ determine every completed right-side
excursion: its exit height, duration and complete inside field after return.
They also determine the complete inside field immediately before the final
entry. Vertically translated request lists give translated full records.
\end{lemma}
\begin{proof}
The first entry sees an entirely white half-plane. Literal dynamics uniquely
determine its first exit and all writes up to that exit, when it exists.
During an outside excursion no inside cell changes. The next entry therefore
sees exactly the retained field produced by the preceding calls. Induction
proves the assertion for all completed excursions. Translation invariance
also identifies histories at horizontally shifted cuts. The final request
records its entry; it need not return.
\end{proof}
All entry heights have the same parity. Put $m=\min_i y_i$ and encode the
record by the integer word
\[
 \rho_c=((y_0-m)/2,\ldots,(y_j-m)/2).
\]
Two genuine records agree up to translation if and only if these request codes
agree: restriction to the entries gives one direction, and
Lemma~\ref{cross:lem:requests} gives the other. Thus equal request codes at two
eligible cuts yield the certificate of Lemma~\ref{cross:lem:records}.

\begin{proposition}[Request capacity]\label{cross:prop:requests}
For $H\ge1$, put $q=\lfloor H/2\rfloor+1$. The number of genuine normalized
records of length $2j+1$ and span at most $H$ is at most
\begin{equation}
 b_H(j)=q^{j+1}-(q-1)^{j+1}.\label{cross:eq:request-capacity}
\end{equation}
For $H=0$, only the one-letter signature exists.
\end{proposition}
\begin{proof}
A request code has $j+1$ letters in $\{0,\ldots,q-1\}$ and contains zero.
Subtracting the words omitting zero gives \eqref{cross:eq:request-capacity}.
Request injectivity makes this an upper bound for genuine signatures. At
height zero a completed return would violate the opposite parities of entry
and exit heights.
\end{proof}
For example, at $H=10$ and $j=5$ the parity capacity is $169{,}643{,}750$,
whereas the request capacity is $31{,}031$. These count combinatorial classes;
a request code contributes an available object even when it has no returning,
span-admissible physical response.

Write $A_H^*(r)=\sum_{j=0}^r b_H(j)$, with $A_H^*(-1)=0$ and
$b_0(j)=\mathbf1_{j=0}$. Let $\mathcal P_H^*(N)$ be the minimum total length
of $N$ distinct objects, with $b_H(j)$ objects of length $2j+1$.
If $r$ is the first index with $A_H^*(r)\ge N$, shortest-layer filling gives
\begin{equation}
 \mathcal P_H^*(N)=\sum_{j<r}(2j+1)b_H(j)
 +(2r+1)\bigl(N-A_H^*(r-1)\bigr).\label{cross:eq:packing}
\end{equation}
Set $\mathcal P_H^*(0)=0$ and $\mathcal P_0^*(N)=\infty$ for $N>1$.
At height one, $b_1(j)=1$, so $\mathcal P_1^*(N)=N^2$.
For $q\ge2$, the cumulative capacity is the integer geometric sum
\[
 A_H^*(r)=\frac{q(q^{r+1}-1)}{q-1}
 -(q-1)\sum_{j=0}^r(q-1)^j.
\]

\begin{theorem}[Request-work dichotomy]\label{cross:thm:work}
Let $N=(p_T.x-a_0)_+$. For $N>0$ let $H$ be the maximum span of the crossing
records at the $N$ cuts $a_0,\ldots,p_T.x-1$. Either a permanent translation
certificate occurs in this finite prefix or
\[
 \mathcal P_H^*(N)\le\lceil T/2\rceil.
\]
\end{theorem}
\begin{proof}
Without a certificate, the $N$ request codes are distinct. Request capacity
therefore bounds their summed full-record lengths below by
\eqref{cross:eq:packing}. Each listed crossing is a different horizontal update,
and horizontal and vertical updates alternate.
\end{proof}

\begin{corollary}[Entropy form]\label{cross:cor:entropy}
In the second alternative, when $H\ge2$ and $N>0$, let
$\overline J=\lfloor(\lceil T/2\rceil-N)/2\rfloor$. Then
\begin{equation}
 \log(N!)\le\overline J\log q
       +N\log\bigl(2(1+\overline J/N)\bigr).
 \label{cross:eq:entropy}
\end{equation}
\end{corollary}
\begin{proof}
The difference-of-powers identity gives $b_H(j)\le(j+1)q^j$, and hence
$A_H^*(r)\le2(r+1)q^r$ for $q\ge2$. Order the pair counts of the $N$ actual
records as $j_1\le\cdots\le j_N$ and put $J=\sum_i j_i$. Since
$i\le A_H^*(j_i)$, summing logarithms and applying concavity gives
\[
 \log(N!)\le J\log q+N\log\bigl(2(1+J/N)\bigr).
\]
Their total length is $2J+N\le\lceil T/2\rceil$, so $J\le\overline J$.
The right-hand side increases with $J$.
\end{proof}

\begin{corollary}[Speed--spread tradeoff]\label{cross:cor:speed}
For an orbit which is not $\ETC$, suppose $T_i\to\infty$ and
$N(T_i)/T_i\to v>0$. Then $v<1/2$ and its record spans satisfy
\[
 \liminf_i\frac{\log(H(T_i)+2)}{\log T_i}\ge\frac{4v}{1-2v}.
\]
\end{corollary}
\begin{proof}
The cases $H\le1$ give at least quadratic work or exclude $N>1$, and are
therefore eventually absent. Also $H\le\lceil T/2\rceil$, since height
variation cannot exceed the number of vertical moves. The work inequality
already gives $v\le1/2$. If $v=1/2$, then $\overline J=o(T)$, and the
right-hand side of \eqref{cross:eq:entropy} is $o(T\log T)$, contradicting
$\log(N!)\sim vT\log T$. For $v<1/2$, divide by $T\log T$ and use
$\overline J/T\to(1-2v)/4$ and $\log N/\log T\to1$.
\end{proof}
Half-turns and quarter-turns give the other directions. For example, speed
$v=1/52$ requires exponent at least $2/25$ on a nontranslating orbit.
Subpolynomial crossing-height spread is consequently incompatible with
positive horizontal upper speed. Together with $H=O(T)$ the bound also gives
$v\le1/6$. No bound on repeated visits is assumed.

\subsection{A numerical finite-profile deadline}
For $H\ge0$, $C\ge1$, and a complete source, define
\begin{align*}
 R&=\max\bigl(\{0\}\cup\{|q.x-p_0.x|:q\in B_0\}\bigr),\\
 L&=H+1,\qquad M(H,C)=A_H^*(LC-1),\\
 D(R,H,C)&=\bigl[2(R+M(H,C))+1\bigr]LC.
\end{align*}
\begin{theorem}[Profile deadline]\label{cross:thm:profile}
Suppose a finite prefix through $T$ has vertical span at most $H$ in every
visited column and at most $C$ updates of every site. Either it contains a
permanent translation certificate or $T\le D(R,H,C)$.
\end{theorem}
\begin{proof}
In an eligible record, east crossings update their departure sites in its left
column. There are at most $LC$ such updates. Since the record starts and ends
east, its odd length is at most $2LC-1$. Every such record therefore lies in a
class of size $M(H,C)$. Reaching a point with
$x>p_0.x+R+M(H,C)$ would give more than $M(H,C)$ cuts beyond the complete initial
support at that subprefix endpoint, forcing a duplicate record. A half-turn
gives the analogous left bound. If no certificate occurs, every head position
has $|x-p_0.x|\le R+M(H,C)$. There are at most $2(R+M)+1$ visited columns,
each with at most $L$ updated sites and at most $C$ updates per site.
\end{proof}
The final unupdated head is not charged as an update. At a final east crossing,
its departure is still in the counted left column; all west crossings of an
eligible record are earlier. Black cells outside the corridor remain present.
Their coordinates matter through the initial frontier, not through an assumed
future-active quotient. The bound is linear in numerical $R$ at fixed $H,C$,
but is not a polynomial bit-complexity claim.

\subsection{Local visit bounds and slowly widening trajectories}
Let $n(v)$ count updates in a finite segment, excluding the terminal position.
For either adjacent pair $u=v+(\pm1,0)$ and $z=v+(0,\pm1)$,
\begin{equation}
 n(v)\le2n(u)+2n(z)+5.\label{cross:eq:corner}
\end{equation}
For completeness, the incoming axis at $v$ is fixed. A visit is exceptional if
it is the initial update, arrives from the selected incoming neighbour, or
leaves to the selected outgoing neighbour. On every other visit both directions
are fixed, hence the turn and read color are fixed. Alternation of colors forces
at least $\lfloor n(v)/2\rfloor$ exceptions. There are at most $n(u)+n(z)+2$,
including the initial and terminal allowances, proving \eqref{cross:eq:corner}.

If all head positions of a prefix lie in $a\le\alpha x+\beta y\le b$ with
$\min(|\alpha|,|\beta|)=1$, choose the adjacent pair increasing the linear form.
Both children increase it by at least one. The analogous pair decreases it.
Iterating toward the nearer boundary, with $F(-1)=0$ and $F(d)=4F(d-1)+5$, gives
\[
 n(v)\le C_h:=\frac{5(4^{\lfloor h/2\rfloor+1}-1)}3,
 \qquad h=\lfloor b-a\rfloor.
\]
Each column contains at most $\lfloor(b-a)/|\beta|\rfloor+1\le h+1$ lattice
sites. Therefore Theorem~\ref{cross:thm:profile} proves the explicit alternative
\begin{equation}
 \text{a certificate occurs, or}\qquad
 T\le D_{\mathrm{obl}}(R,h):=D(R,h,C_h).\label{cross:eq:deadline}
\end{equation}
No infinite confinement is assumed in \eqref{cross:eq:deadline}.

\begin{proof}[Proof of Theorem~\ref{cross:thm:gap}]
For $h\ge2$, $C_h\le(20/3)2^h$ and $q\le h+1$. Bounding $A_h^*$ by
$2(r+1)q^r$ at $r=(h+1)C_h-1$ yields
\[
 \log D_{\mathrm{obl}}(R,h)
 =O_R\bigl((h+1)2^h\log(h+2)\bigr).
\]
On an orbit which is not $\ETC$, no finite certificate can occur. Thus for every
$T$, with $h=\lfloor w_{\alpha,\beta}(T)\rfloor$,
\[
 \log\log T\le h\log2+\log(h+1)+\log\log(h+2)+O_R(1).
\]
This first forces $h\to\infty$ and then gives the claimed liminf. All constants
are independent of the normalized normal. For an infimum choose a normal with
width less than one above the infimum; the bounded difference does not change
the asymptotic conclusion.
\end{proof}

For example, an eventual tube
\[
 |y-\theta x|\le g(|x|),\qquad \theta\ne0,
 \qquad g(r)=o\bigl(\log\log(r+e^e)\bigr)
\]
forces $\ETC$: unit speed makes its prefix normal span $o(\log\log T)$; a finite
initial part contributes only a constant. For irrational $\theta$ this is
impossible, because a nonzero integer translating drift has a nonzero linear
deviation from that slope. The width $g$ is allowed to tend to infinity.

There is a corresponding nonstraight extension. Let $f:\Z\to\Z$ be
nondecreasing with $f(x+L)\ge f(x)+1$ for some fixed $L\ge1$. A north/west path
of length $L(H+1)$ leaves the corridor $f(x)\le y\le f(x)+H$ above its top.
Expanding \eqref{cross:eq:corner} along this binary tree gives
\[
 C_{L,H}=\frac{5(4^{L(H+1)}-1)}3.
\]
The profile deadline then implies, for any nontranslating orbit whose prefix
fits that corridor with a possibly increasing $H(T)$,
\[
 \liminf_{T\to\infty}\frac{H(T)}{\log\log T}\ge\frac1{L\log4}.
\]
The bound $H(T)$ covers the entire prefix, not a window from which old deviations
have been deleted.

\subsection{Nonconstant cyclic turning rules}
The geometric information argument also applies to generalized cyclic ants.
Let a rule $\omega\in\{R,L\}^m$ contain both letters, let a read increment its
cell state modulo $m$, and let only finitely many initial states differ from
zero. Put
\[
 a=\min(\#R(\omega),\#L(\omega)),\qquad \kappa=m/a.
\]
For a cyclic prefix at phase $s$, let $n_c(s,r)$ count turn $c$. Define
\[
 b_\omega=\max_{s,c,\,0\le r<m}\{r-\kappa n_c(s,r)\}.
\]
Every full $m$-block has at least $a$ copies of either turn, so
$n\le\kappa n_c(s,n)+b_\omega$ for every prefix length. In the corner argument,
all nonexceptional visits require one fixed turn. Visits with the opposite turn
are exceptional, giving
\[
 n(v)\le\kappa(n(u)+n(z))+2\kappa+b_\omega.
\]
The rest of the proof is unchanged, replacing $C_h$ by
\[
 C_\omega(h)=\left\lceil
 (2\kappa+b_\omega)
 \frac{(2\kappa)^{\lfloor h/2\rfloor+1}-1}{2\kappa-1}
 \right\rceil.
\]

\begin{theorem}[Frequency-dependent width gap]
Every non-eventually-translating finite-support orbit for the fixed rule
$\omega$ satisfies
\[
 \liminf_{T\to\infty}\frac{w_{\alpha,\beta}(T)}{\log\log T}
 \ge\frac2{\log(2\kappa)}.
\]
\end{theorem}
\begin{proof}
Heading-axis parity and crossing parity still alternate. Complete
half-plane coupling and request reconstruction remain deterministic with $m$ colors. The local visit
bound grows as $O_\omega((2\kappa)^{h/2})$. Substitution into the same finite
profile deadline and taking two logarithms proves the claim.
\end{proof}
For classical $RL$, $\kappa=2$, $b_\omega=1$, recovering the exact earlier
local-time constant. Every balanced cyclic rule has the same asymptotic
coefficient $1/\log2$, though its finite constants may differ.

\subsection{Effectivity and reproducibility}
For a supplied rational non-axis normal and integer width bound, an implemented
algorithm follows complete states, checks the running width, and compares all
crossing records at new left/right horizontal records beyond the full initial
support. It returns a first observed width violation or a permanent translating
certificate. Without a user cap, \eqref{cross:eq:deadline} proves termination. It need
not construct the large deadline integer during execution. A capped run returns
\texttt{UNKNOWN}, never a negative dynamical conclusion.

The certificate checker independently replays the finite source, compares the
full half-plane fields, and checks the finite no-exit guard. A separate phase-ray
parity check tests all initial black hits, all self-write breakpoints, and the
final infinite interval. An unseen forward defect is therefore not ignored.
A translating answer does not assert permanent satisfaction of the supplied
width bound; a non-tangent translator may later leave it.

The Python simulator and an independently written C++ vector simulator agree
on every path position and complete endpoint field of 100 supplied inputs,
totalling 113,152 updates per implementation. Further checks exhaust the
stated small parity-word classes, verify integer packing and entropy inequalities,
check finite corner bounds, and reject malformed or unguarded certificates.
Four actual standard-tail inputs in different rotations return 104-step
certificates. These are implementation checks, not proofs that all tails are
standard. The finite cyclic-rule balance calculation is checked for every
nonconstant rule word of lengths two through eight. The general-rule classifier
is not implemented; its extension is analytic.

The request-capacity module separately checks 78 complete inputs and 86,647
updates per literal implementation. A persistent half-plane replay compares
186 complete residual fields; the controls also compare request-code and
full-signature equality, exact integer packing, and damaged histories.
It uses the same hash-pinned baseline simulator and C++ oracle as the preceding
module, together with a separate right-excursion implementation.

The reproduction entry point is
\begin{quote}\ttfamily
python3 memory\_feedback/reproduce.py --crossing-only
\end{quote}
using Python's standard library and a C++17 compiler. The default replay works
in a temporary directory and compares an immutable report. The universal
claims are the displayed analytic arguments, not extrapolations of the tests.



\section{Reproduction}\label{reproduction-and-perspective}

Run \texttt{python3\ memory\_feedback/reproduce.py} to replay the eight
scoped modules: first-event prediction, exact XOR continuation, event
geometry, remote filtering, birth-filtered feedback, the twenty-cell circuit, crossing information, and request reconstruction. The checks
compare full finite fields, reconstruct integer and rational
inequalities, and verify permanent periodic tails at all parity
breakpoints. Parameter families are covered by the affine and geometric
proofs above; finite controls test their numerical premises.

\section{Limitations}

Request capacity counts possible incoming codes; it is an upper bound on genuine
crossing histories, not a physical realization count or a stock fixed by the
initial state. Request injectivity proves uniqueness of completed excursions;
it does not decide whether an arbitrary proposed excursion returns. The work
bound concerns distinct records, without assuming statistically independent
input information.
\label{limitations}
The width estimate uses whole-prefix normal-coordinate width with both normal
components nonzero. The exact combinatorial capacity does not establish physical
sharpness of the asymptotic width coefficient. A translating certificate allows
general periods and drifts. The extension to nonconstant cyclic turning rules is
analytic; a generalized-rule classifier is not implemented.


The polynomial bounds fix \(D\), and the spatially filtered query also
fixes \(K\) in the normalized input frame. The factor \(f(D)\) may be
doubly exponential, and the finite-region escape constant can be enormous.
The event count \(n\) is measured numerically, not by \(\log n\);
the bounds do not give a polynomial ordinary-step simulation. Arbitrary
compressed input strides retain totality but need not obey the bounded
residue periods of histories generated from explicit sparse inputs.
For fixed \(K\), the exterior deadline has a leading coefficient
independent of \(K\); its linear coefficient and asymptotic onset depend
on \(b_K\). This gives no uniform estimate for regions growing with time.

The total queries decide a next event or a finite event prefix. They do
not decide whether an infinite event chain occurs, nor whether every
finite seed eventually reaches a standard highway. The directional tube
bounds control ranges rather than limiting speeds or angular convergence.
Remote feedback is necessary on the stated non-eventually-translating
orbits; the construction supplies a family of increasingly long lifetimes
with different finite inputs. Each member reaches the standard highway,
so autonomous infinite renewal from one fixed seed remains open.

\clearpage
\section{Perspective}

The predictor supplies a total query, while the XOR representation
supplies its reusable physical output. The tube and escape arguments
make fixed-region filtering total and force nonperiodic behaviour to
renew memory remotely. The construction shows that mass and event count
alone cannot replace spatial information. The crossing argument supplies a
complementary requirement: persistent nontranslation forces every normalized
oblique width to grow at least double-logarithmically with elapsed time.
Complete-field coupling connects both parts: it makes event continuation exact
and turns repeated guarded crossing records into permanent translation.

\section{References}\label{references}

\begingroup\small
\begin{enumerate}
\def\labelenumi{\arabic{enumi}.}
\item
  Anahí Gajardo. \emph{A symbolic projection of Langton's Ant (Extended
  Abstract)}. DMTCS Proceedings AB, 57--68 (2003).
  \href{https://doi.org/10.46298/dmtcs.2312}{doi:10.46298/dmtcs.2312}.
\item
  Anahí Gajardo, Victor H. Lutfalla and Michaël Rao. \emph{Ants on the
  highway}. Natural Computing 24, 497--509 (2025).
  \href{https://doi.org/10.1007/s11047-025-10018-9}{doi:10.1007/s11047-025-10018-9}.
  Author version: \href{https://arxiv.org/abs/2409.10124v1}{arXiv:2409.10124v1}.
\item
  Atharva Jillhewar. \emph{Finite-Support Periodic Highways of Langton's
  Ant: Necessary Conditions, Transverse Exclusions, and Exact Search}
  (2026).
  \href{https://doi.org/10.21203/rs.3.rs-10516269/v1}{doi:10.21203/rs.3.rs-10516269/v1}.
  Theorem 3.1 and Remark 3.2 supply full-drift visit chronology and seed
  reconstruction.
\item
  Hao Ke. \emph{An Ant on the Clock: Exponential Escape Bounds for
  Langton's Ant}, version 0.4.0 (2026).
  \href{https://doi.org/10.5281/zenodo.23402379}{doi:10.5281/zenodo.23402379}.
\item
  Hao Ke. \emph{Four Black Cells Force an Outgoing Highway: A Sharp
  Half-Plane Threshold for Langton's Ant}, version 0.1.0 (2026).
  \newline\href{https://doi.org/10.5281/zenodo.23402414}{doi:10.5281/zenodo.23402414}.
\item
  Anahí Gajardo, Andrés Moreira and Eric Goles. \emph{Complexity of
  Langton's ant}. Discrete Applied Mathematics 117, 41--50 (2002).
  \href{https://doi.org/10.1016/S0166-218X(00)00334-6}{doi:10.1016/S0166-218X(00)00334-6}.
\item F. C. Hennie. One-Tape, Off-Line Turing Machine Computations.
\emph{Information and Control} 8 (1965), 553--578.
\href{https://doi.org/10.1016/S0019-9958(65)90399-2}{doi:10.1016/S0019-9958(65)90399-2}.
\item D. Ranjan, R. Chang and J. Hartmanis. \emph{Space Bounded Computations:
Review and New Separation Results}. Cornell University report,
20 December 1989, Section 2.
\url{https://userpages.cs.umbc.edu/chang/papers/loglog/loglog.pdf}.
\end{enumerate}
\endgroup

\end{document}
